%==================================================
% Plantilla elaborado por Dr. David Pareja Quispe
% e-mail: dparejaq@unmsm.edu.pe
% No eliminar los créditos del Autor
%================================================
\documentclass[twoside,onehalfspace]{./Portada/unmsmtesis}

%=========================================================
% ADICIONAR LINEAS ABAJO OTROS PAQUETES SI ES NECESARIO
%=========================================================
\graphicspath{{Figuras/}}%Ruta de las figuras
\usepackage{blindtext}

%==========================================
%    Fuente
%==========================================
\newcommand*{\captionsource}[2]{%
  \caption[{#1}]{#1}\smallskip\par
  \textbf{Fuente:} #2\par}


%===========================================
% Segunda opción para la fuente
%==========================================
%\usepackage{copyrightbox} 
%\makeatletter
%\renewcommand{\CRB@setcopyrightfont}{%
%\footnotesize  
%\color{black} %Comentar si desea en el color por defecto
%}

%%%%%%%%%%%%%%%%%%%%%%%%%%%%%%%%%%%%%%%%%%%%%%%%%%%%%%%%%%%%%
%%%%%%%%%%%%%% AQUÍ INICIA EL DOCUMENTO %%%%%%%%%%%%%%%%%%%%%
%%%%%%%%%%%%%%%%%%%%%%%%%%%%%%%%%%%%%%%%%%%%%%%%%%%%%%%%%%%%%
%\setlength {\marginparwidth}{2cm} 
\begin{document}

%%%%%%%%%%%%%%%%%%%%%%%%%%%%%%%%%%%%%%%%%%
% CARGANDO LA PORTADA DEL DOCUMENTO
\begin{titlepage}
	\centering
	\vspace{-20pt}
    \hrule {\color{purple} \rule{\linewidth}{0.5mm}} \hrule 
    %\vspace{0.05cm}
    %\hrule
	\vspace{0.8cm}
	\textsc{\Large \textbf{Universidad Nacional Mayor San Marcos}\\}
	\vspace{0.5cm}
	\textbf{\normalsize Universidad del Perú, DECANA DE AMÉRICA\\}
    \vspace*{0.8cm}
    \hrule  {\color{purple} \rule{\linewidth}{0.5mm}} \hrule
	\vspace{1.2cm}
	\textsc{\Large \bfseries Facultad de Ciencias Físicas \\}
    \vspace{0.4cm}
	\textbf{\large \bfseries Escuela Profesional de Física \\}
    %\vspace{0.05cm}
    %\hrule
    \vspace*{1.0 cm}
    \includegraphics[scale = 0.10]{Portada/logoUNMSM.png}\\[1.0 cm]	% UDG Logo
    \vspace{1.0cm}
    %%%%%%%%%%%%%%%%%%%%%%%%%%%%%%%%%%%%%%%%%%%%%%%%%%%
    {\large \bfseries Análisis perturbativo de la esfera de fotones y detectabilidad de la sombra de un agujero negro en el espaciotiempo de Kerr-Newman con carga débil}\\
    %%%%%%%%%%%%%%%%%%%%%%%%%%%%%%%%%%%%%%%%%%%%%%%%%%%
    \vspace{1.5cm}
    {\normalsize \mdseries Henry Paolo Alfaro Sotil \\
    } 
%==========================================
%              LICENCIATURA
%RETIRAR EL COMENTARIO SEGUN CORRESPONDA
%==========================================
    \vspace{1.5cm}
    {\normalsize 
    Tesis para optar el Título Profesional de Licenciado en Física, presentado a la Escuela Profesional de Física, asesorado por el Dr. David Pareja Quispe, aprobado el XX de mes de 202X.}

%==========================================
%                 MAESTRIA
%RETIRAR EL COMENTARIO SEGUN CORRESPONDA
%=========================================
%    \vspace{1.5cm}
%    {\normalsize 
%    Tesis para optar el grado de Magíster en Geofísica, presentado a la Escuela Profesional de Ciencias Físicas, asesorado por el Dr. David Pareja Quispe, aprobado el XX de octubre del 202X.}

%===========================================
%                 DOCTORADO
%RETIRAR EL COMENTARIO SEGUN CORRESPONDA
%===========================================
%    \vspace{1.5cm}
%    {\normalsize 
%    Tesis para optar el grado de Doctor en Geofísica, presentado a la Escuela Profesional de Ciencias Físicas, asesorado por el Dr. David Pareja Quispe, aprobado el XX de octubre del 202X.}

    %%%%%%%%%%%%%%%%%%%%%%%%%%%%%%%%%%%%%%%%%%%%%%%%%%%
    %\raggedright
    %{\normalsize Asesor: Dr. David Pareja Quispe \\
    %}  
    \center
    \vspace{4\baselineskip}
    {\large  Lima - Perú}\\
%    \vspace{0.3cm}
\medskip
    {\large \textbf{\the\year}} \\\vfill %AÑO DE PUBLICACIÓN
%%%%%%%%%%%%%%%%%%%%%%%%%%%%%%%%%%%%%%%%%%%%%%%%%%%    
\end{titlepage}
%%%%%%%%%%%%%%%%%%%%%%%%%%%%%%%%%%%%%%%%%%
\thispagestyle{empty}\null\newpage
%%%%%%%%%%%%%%%%%%%%%%%%%%%%%%%%%%%%%%%%%%
% REIMPRESIÓN DE LA PORTADA NUEVAMENTE
%\begin{titlepage}
	\centering
	\vspace{-20pt}
    \hrule {\color{purple} \rule{\linewidth}{0.5mm}} \hrule 
    %\vspace{0.05cm}
    %\hrule
	\vspace{0.8cm}
	\textsc{\Large \textbf{Universidad Nacional Mayor San Marcos}\\}
	\vspace{0.5cm}
	\textbf{\normalsize Universidad del Perú, DECANA DE AMÉRICA\\}
    \vspace*{0.8cm}
    \hrule  {\color{purple} \rule{\linewidth}{0.5mm}} \hrule
	\vspace{1.2cm}
	\textsc{\Large \bfseries Facultad de Ciencias Físicas \\}
    \vspace{0.4cm}
	\textbf{\large \bfseries Escuela Profesional de Física \\}
    %\vspace{0.05cm}
    %\hrule
    \vspace*{1.0 cm}
    \includegraphics[scale = 0.10]{Portada/logoUNMSM.png}\\[1.0 cm]	% UDG Logo
    \vspace{1.0cm}
    %%%%%%%%%%%%%%%%%%%%%%%%%%%%%%%%%%%%%%%%%%%%%%%%%%%
    {\large \bfseries Análisis perturbativo de la esfera de fotones y detectabilidad de la sombra de un agujero negro en el espaciotiempo de Kerr-Newman con carga débil}\\
    %%%%%%%%%%%%%%%%%%%%%%%%%%%%%%%%%%%%%%%%%%%%%%%%%%%
    \vspace{1.5cm}
    {\normalsize \mdseries Henry Paolo Alfaro Sotil \\
    } 
%==========================================
%              LICENCIATURA
%RETIRAR EL COMENTARIO SEGUN CORRESPONDA
%==========================================
    \vspace{1.5cm}
    {\normalsize 
    Tesis para optar el Título Profesional de Licenciado en Física, presentado a la Escuela Profesional de Física, asesorado por el Dr. David Pareja Quispe, aprobado el XX de mes de 202X.}

%==========================================
%                 MAESTRIA
%RETIRAR EL COMENTARIO SEGUN CORRESPONDA
%=========================================
%    \vspace{1.5cm}
%    {\normalsize 
%    Tesis para optar el grado de Magíster en Geofísica, presentado a la Escuela Profesional de Ciencias Físicas, asesorado por el Dr. David Pareja Quispe, aprobado el XX de octubre del 202X.}

%===========================================
%                 DOCTORADO
%RETIRAR EL COMENTARIO SEGUN CORRESPONDA
%===========================================
%    \vspace{1.5cm}
%    {\normalsize 
%    Tesis para optar el grado de Doctor en Geofísica, presentado a la Escuela Profesional de Ciencias Físicas, asesorado por el Dr. David Pareja Quispe, aprobado el XX de octubre del 202X.}

    %%%%%%%%%%%%%%%%%%%%%%%%%%%%%%%%%%%%%%%%%%%%%%%%%%%
    %\raggedright
    %{\normalsize Asesor: Dr. David Pareja Quispe \\
    %}  
    \center
    \vspace{4\baselineskip}
    {\large  Lima - Perú}\\
%    \vspace{0.3cm}
\medskip
    {\large \textbf{\the\year}} \\\vfill %AÑO DE PUBLICACIÓN
%%%%%%%%%%%%%%%%%%%%%%%%%%%%%%%%%%%%%%%%%%%%%%%%%%%    
\end{titlepage}

%================================================
%ESPACIAMIENTO ENTRE LINEAS EN TODO EL DOCUMENTO
%================================================
\onehalfspacing 
\pagenumbering{roman}

%=======================================================================================
%  COMENTE LAS LINEAS (PARTES DEL DOCUMENTO) QUE NO DESEA QUE APAREZCA EN EL ÍNDICE GENERAL
%=======================================================================================
%==========================
%          EPÍGRAFE
%==========================
\chapter*{\hfill{\centering Epígrafe}\hfill}
\addcontentsline{toc}{chapter}{Epígrafe}  %ADICIONAR O RETIRAR el comentario si quiere que aparesca o que no aparesca en la lista del Indice general
\normalsize

\vspace{11cm} % Cambiar este espacio (11cm) de acuerdo a su interes


\justifying{\Large\itshape{``Puedes estudiar y hacer lo que quieras siempre y cuando
    en eso que seas, seas lo mejor de lo mejor''}}   %CAMBIAR EL TEXTO 

\large\raggedleft\textbf{\textit{Henry Alfaro}}  %CAMBIAR EL TEXTO
 %% OPCIONAL
%%%%%%%%%%%%%%%%%%%%%%%%%%%%%%%%%%%%%%%%%%
%               DEDICATORIA
%=========================================
\chapter*{\hfill{\centering Dedicatoria}\hfill}
\addcontentsline{toc}{chapter}{Dedicatoria}  %Retirar comentario si quiere que aparesca en la lista del Indice general
\normalsize
\justify

\vspace{19cm} % Cambiar este espacio (19 cm) de acuerdo en que parte de la hoja quiere mostrar el texto de su dedicatoria

\begin{flushright} 
\itshape\Large 

\textit{A mis padres \textbf{Valeria} y \textbf{Henry}, también a mi tío Carlos y a mi tía Nathaly, a mis abuelos Guillermina y Hugo.} %CAMBIAR EL TEXTO 

\end{flushright}

 %% OPCIONAL
%%%%%%%%%%%%%%%%%%%%%%%%%%%%%
%       AGRADECIMIENTOS
%============================
\chapter*{\hfill{\centering Agradecimientos}\hfill}
\addcontentsline{toc}{chapter}{Agradecimientos}  A Dios, por ser mi refugio y fortaleza en los momentos de crisis, dolor y miedo. Gracias por llenar mi vida de resiliencia para afrontar cada reto, por motivarme a seguir su voluntad con una fe ciega e inquebrantable, y por guiar cada uno de mis pasos a lo largo de este camino.

A mis padres, por su fe inquebrantable en mí y por haberme brindado el regalo más valioso de mi vida: la educación. Su esfuerzo y sacrificio han sido el pilar fundamental de cada uno de mis logros.

A mis abuelos, por haberme acogido en un hogar donde pude crecer, aprender de su sabiduría y dedicarme al estudio sin límites. Un agradecimiento especial a mi abuela, por haber financiado mis estudios escolares y por creer siempre en mi potencial.

A mi tío Carlos, por su apoyo incondicional en cada etapa de mi vida. Gracias por estar a mi lado durante mi trayectoria en el fútbol profesional, por respaldarme en el momento de mi retiro y por impulsarme con convicción a emprender el camino universitario. Su ayuda económica, sus gestos constantes y su respaldo con los materiales de estudio fueron clave para construir este nuevo capítulo.

A mi tía Nathaly, por su generosidad y confianza al costear mi preparación preuniversitaria en cuanto supo de mi vocación por postular a mi alma mater. Su gesto fue el primer paso hacia esta meta académica.

Un agradecimiento muy especial a mi amigo Javier, por haberme mostrado la belleza de la física, enseñándome lo que significa ser un físico y lo fascinante que es estudiar esta carrera. Gracias por ser un mentor en una etapa de mi vida en la que enfrentaba serias carencias académicas, dedicando su tiempo a enseñarme matemática y física, y encendiendo en mí la pasión por este camino.

A mi comunidad Luz de Cristo, por abrirme sus puertas con un amor fraternal y ser un espacio fundamental de sanación y crecimiento. Gracias por haberme brindado un entorno donde aprendí a exponer mis ideas con confianza, por la formación moral y en la fe, y por la valiosa presencia de compañeros universitarios cuyos consejos me ayudaron a sobrellevar la incertidumbre sobre el futuro de mi carrera.

A mis amigos y hermanos de vida, por su cariño incondicional, por cuidar de mí y motivarme a la superación constante. Gracias de corazón por su empatía y solidaridad en aquellos momentos en los que mis recursos económicos eran escasos, cubriendo generosamente mis pasajes, comidas o entradas para que nunca dejara de compartir y disfrutar con ustedes.

Finalmente, a cada uno de mis profesores y entrenadores a lo largo de mi vida, por exigirme al máximo en cada aula, examen y entrenamiento, enseñándome a dar siempre la mejor versión de mí mismo
\vspace{0.5cm}
\justify
\normalsize

%=================================


\newpage %% OPCIONAL
%==========================
%          RESUMEN
%==========================
\chapter*{\hfill{\centering Resumen}\hfill}
\addcontentsline{toc}{chapter}{Resumen}  %Retirar comentario si quiere que aparesca en la lista del Indice general
\vspace{0.5cm}
\justify
\normalsize

El presente trabajo de investigación analiza analítica y perturbativamente la dinámica geodésica de fotones y la estructura de la sombra en el espaciotiempo de Kerr-Newman bajo la aproximación de carga eléctrica débil ($\delta \equiv \frac{Q^2}{4\pi\varepsilon_0 G M^2} \ll 1$), formulado rigurosamente en unidades del Sistema Internacional (SI) y empleando la convención de firma métrica de Minkowski $(+,-,-,-)$. Se deduce paso a paso el potencial radial efectivo $\mathcal{R}(r)$ a partir del formalismo de Hamilton-Jacobi y la función de acción $S$, obteniendo las constantes conservadas de energía $E$, momento angular azimutal $L_z$ y la constante de Carter $\mathcal{K}$. Definiendo los parámetros de impacto reducidos $\Phi \equiv \frac{c L_z}{E}$ y $\mathcal{Q} \equiv \frac{c^2 \mathcal{K}}{E^2}$, y estableciendo las condiciones de órbita fotónica esférica inestable $\mathcal{R}(r_{\mathrm{ph}}) = 0$ y $\left.\frac{d\mathcal{R}}{dr}\right|_{r_{\mathrm{ph}}} = 0$, se extiende la metodología de Kobialko para métricas rotacionales no diagonales a partir del parámetro crítico $b_c^2 = \mathcal{Q}_c + \Phi_c^2$. Se deriva la expansión perturbativa del radio efectivo equivalente $R_{\mathrm{eff}}(\delta)/R_{\mathrm{eff}}(0) = 1 - \frac{1}{6}\delta - \frac{7}{216}\delta^2$, el coeficiente de segundo orden $D = 1/27$ en $R_{\mathrm{eff}}^2$ que aísla la firma unívoca de electrovacío frente a modelos de gravedad modificada, y la razón de desplazamiento del primer nulo de visibilidad $u_1(\delta)/u_1(0) = 1 + \frac{1}{6}\delta + \frac{13}{216}\delta^2$. Asimismo, se demuestra que mientras en el límite estático de Schwarzschild el factor de adelgazamiento de Sneppen ($f_{\mathrm{thin}} = e^{-2\pi} \approx 1/535$) vuelve inobservable el subanillo $n=1$, en Kerr-Newman el momento angular degrada el exponente de Lyapunov ($\gamma_\phi < 1$), ensanchando orbitalmente el primer subanillo de fotones. Finalmente, se integran las geodésicas nulas mediante el algoritmo Runge-Kutta de cuarto orden (RK4) en Python y se contrastan los perfiles teóricos de visibilidad interferométrica $|V(u)|$ en el plano de Fourier con las observaciones del \textit{Event Horizon Telescope} (EHT) para $\text{Sgr A}^*$ (repositorio 2024-D02-01) y $\text{M87}^*$ (repositorio 2024-D01-01). Se encuentra que para $\text{Sgr A}^*$, el modelo de Kerr-Newman con $\delta \approx 0.290$ desplaza el primer nulo de visibilidad en un $+5.42\%$, reduciendo la bondad de ajuste en $\Delta\chi^2_\nu = -4.32$ ($\Delta\chi^2 = -40\,166$) frente a Schwarzschild; mientras que en $\text{M87}^*$, el modelo simétrico 1D empeora ($\Delta\chi^2_\nu = +10.74$) debido a la contaminación por la emisión extendida del chorro relativista. Se concluye que el ajuste en $\text{Sgr A}^*$ actúa como un compensador de escala geométrica en el modelo simétrico 1D y no constituye una detección física directa de carga, dada la dominancia de las incertidumbres sistemáticas actuales en la razón masa-distancia ($\sim 11.8\%$).

%=================================
% Palabras claves
%=================================
\vspace{0.5cm}
\raggedright 
{\bf \textit{Palabras clave}:} Agujero negro de Kerr-Newman, Carga débil, Esfera de fotones, Radio efectivo de la sombra, Constante de Carter, Interferometría VLBI, Event Horizon Telescope, Exponente de Lyapunov.
 %% OBLIGATORIO 
%==========================
%          ABSTRACT
%==========================
\chapter*{\hfill{\centering Abstract}\hfill}
\addcontentsline{toc}{chapter}{Abstract}  %Retirar comentario si quiere que aparesca en la lista del Indice general
\vspace{0.5cm}
\justify
\normalsize

This research work presents an analytical and perturbative analysis of photon geodesic dynamics and shadow structure in the Kerr-Newman spacetime under the weak electric charge approximation ($\delta \equiv \frac{Q^2}{4\pi\varepsilon_0 G M^2} \ll 1$), rigorously formulated in International System (SI) units and employing the Minkowski metric signature convention $(+,-,-,-)$. The effective radial potential $\mathcal{R}(r)$ is derived step by step using the Hamilton-Jacobi formalism and the action function $S$, yielding the conserved constants of motion: energy $E$, azimuthal angular momentum $L_z$, and Carter's constant $\mathcal{K}$. Defining the reduced impact parameters $\Phi \equiv \frac{c L_z}{E}$ and $\mathcal{Q} \equiv \frac{c^2 \mathcal{K}}{E^2}$, and imposing the unstable spherical photon orbit conditions $\mathcal{R}(r_{\mathrm{ph}}) = 0$ and $\left.\frac{d\mathcal{R}}{dr}\right|_{r_{\mathrm{ph}}} = 0$, Kobialko's methodology is extended to non-diagonal rotational metrics via the critical parameter $b_c^2 = \mathcal{Q}_c + \Phi_c^2$. We derive the perturbative expansion of the equivalent effective radius $R_{\mathrm{eff}}(\delta)/R_{\mathrm{eff}}(0) = 1 - \frac{1}{6}\delta - \frac{7}{216}\delta^2$, the second-order coefficient $D = 1/27$ in $R_{\mathrm{eff}}^2$ isolating the unique electrovacuum charge signature from modified gravity models, and the displacement ratio of the first visibility null $u_1(\delta)/u_1(0) = 1 + \frac{1}{6}\delta + \frac{13}{216}\delta^2$. Furthermore, it is shown that while in the static Schwarzschild limit Sneppen's thinning factor ($f_{\mathrm{thin}} = e^{-2\pi} \approx 1/535$) renders the $n=1$ subring unobservable, in Kerr-Newman the angular momentum degrades the Lyapunov exponent ($\gamma_\phi < 1$), orbitally widening the first photon subring. Finally, null geodesics are integrated using the fourth-order Runge-Kutta (RK4) algorithm in Python, and the theoretical interferometric visibility profiles $|V(u)|$ in the Fourier domain are contrasted with Event Horizon Telescope (EHT) observations for $\text{Sgr A}^*$ (2024-D02-01 data release) and $\text{M87}^*$ (2024-D01-01 data release). For $\text{Sgr A}^*$, the Kerr-Newman model with $\delta \approx 0.290$ shifts the first visibility null by $+5.42\%$, improving the goodness-of-fit by $\Delta\chi^2_\nu = -4.32$ ($\Delta\chi^2 = -40\,166$) compared to Schwarzschild; whereas for $\text{M87}^*$, the 1D symmetric model degrades ($\Delta\chi^2_\nu = +10.74$) due to contamination from extended relativistic jet emission. It is concluded that the fit in $\text{Sgr A}^*$ acts as a geometric scale compensator within the 1D symmetric model rather than a direct physical charge detection, given the dominance of current systematic uncertainties in the mass-to-distance ratio ($\sim 11.8\%$).


%=================================
% Keywords
%=================================
\vspace{0.5cm}
\raggedright 
{\bf \textit{Keywords}:} Kerr-Newman black hole, Weak electric charge, Photon sphere, Effective shadow radius, Carter's constant, VLBI Interferometry, Event Horizon Telescope, Lyapunov exponent.

 %% OBLIGATORIO

%%%%%%%%%%%%%%%%%%%%%%%%%%%%%%%%%%%%%%%%%%
%             ÍNDICE GENERAL
%%%%%%%%%%%%%%%%%%%%%%%%%%%%%%%%%%%%%%%%%%
\tableofcontents
%\pagebreak
\newpage
%=================================================
%%%%%%%%%%%%%%%%%%% LISTA DE FIGURAS  %%%%%%%%%%%%
%=================================================
\listoffigures
\addcontentsline{toc}{chapter}{\listfigurename} % OBLIGATORIO SI HAY MAS DE 3 FIGURAS
\null\newpage
%\thispagestyle{empty}\null\newpage
%=========================================
%%%%%%%%%%%%%%%%LISTA DE TABLAS%%%%%%%%%%%
%=========================================
\listoftables
\addcontentsline{toc}{chapter}{\listtablename} % OBLIGATORIO SI HAY MAS DE 3 TABLAS
\null\newpage
%\thispagestyle{empty}\null\newpage

%=========================================
% LISTA DE SIMBOLOS Y SIGLAS ES OPCIONAL
%=========================================
%%%%%%%%%%%%%%%%%%%%%%%%%%%%%%%%%%%%%%%%
%       LISTA DE SÍMBOLOS
%%%%%%%%%%%%%%%%%%%%%%%%%%%%%%%%%%%%%%%%
\chapter*{\hfill{\centering Lista de símbolos}\hfill}
\addcontentsline{toc}{chapter}{Lista de símbolos}
\vspace{0.5cm}
\normalsize

\noindent\textbf{Constantes fundamentales y magnitudes métricas}
\vspace{0.2cm}

\begin{tabular}{ll}
$c$ & --$\quad$ Velocidad de la luz en el vacío ($\text{m}\cdot\text{s}^{-1}$)\\
$G$ & --$\quad$ Constante de gravitación universal ($\text{m}^3\cdot\text{kg}^{-1}\cdot\text{s}^{-2}$)\\
$\varepsilon_0$ & --$\quad$ Permitividad eléctrica del vacío ($\text{F}\cdot\text{m}^{-1}$)\\
$\mu_0$ & --$\quad$ Permeabilidad magnética del vacío ($\text{H}\cdot\text{m}^{-1}$)\\
$M$ & --$\quad$ Masa del agujero negro ($\text{kg}$)\\
$Q$ & --$\quad$ Carga eléctrica neta del agujero negro ($\text{C}$)\\
$J$ & --$\quad$ Momento angular del agujero negro ($\text{kg}\cdot\text{m}^2\cdot\text{s}^{-1}$)\\
$a \equiv \frac{J}{Mc}$ & --$\quad$ Parámetro de espín por unidad de masa ($\text{m}$)\\
$\delta \equiv \frac{Q^2}{4\pi\varepsilon_0 G M^2}$ & --$\quad$ Parámetro adimensional de carga eléctrica débil ($\delta \ll 1$)\\
$r_g \equiv \frac{GM}{c^2}$ & --$\quad$ Radio gravitacional ($\text{m}$)\\
$\theta_g \equiv \frac{r_g}{D}$ & --$\quad$ Escala angular gravitacional ($\text{rad}$ o $\mu\text{as}$)\\
\end{tabular}

\vspace{0.5cm}
\noindent\textbf{Tensores y funciones métricas en $(+,-,-,-)$}
\vspace{0.2cm}

\begin{tabular}{ll}
$g_{\mu\nu}$ & --$\quad$ Tensor métrico del espaciotiempo con firma $(+,-,-,-)$\\
$g \equiv \det(g_{\mu\nu})$ & --$\quad$ Determinante del tensor métrico ($g < 0$)\\
$R_{\mu\nu}$ & --$\quad$ Tensor de curvatura de Ricci\\
$R$ & --$\quad$ Escalar de curvatura de Ricci ($R = g^{\mu\nu} R_{\mu\nu}$)\\
$F_{\alpha\beta}$ & --$\quad$ Tensor de fuerza del campo electromagnético\\
$A_\mu$ & --$\quad$ Cuadripotencial electromagnético ($\text{V}\cdot\text{s}\cdot\text{m}^{-1}$)\\
$T_{\mu\nu}^{(\mathrm{EM})}$ & --$\quad$ Tensor de energía-impulso electromagnético ($\text{J}\cdot\text{m}^{-3}$)\\
$\Sigma(r,\theta)$ & --$\quad$ Función métrica auxiliar de Boyer-Lindquist: $r^2 + a^2 \cos^2\theta$ ($\text{m}^2$)\\
$\Delta(r)$ & --$\quad$ Función de horizonte de Kerr-Newman ($\text{m}^2$)\\
\end{tabular}
\newpage
\vspace{0.5cm}
\noindent\textbf{Dinámica de fotones, constantes de movimiento y teoría perturbativa}
\vspace{0.2cm}

\begin{tabular}{ll}
$S$ & --$\quad$ Función de acción principal de Hamilton-Jacobi ($\text{J}\cdot\text{s}$)\\
$E$ & --$\quad$ Energía del fotón conservada por simetría temporal ($\text{J}$)\\
$L_z$ & --$\quad$ Momento angular azimutal conservado por simetría axial ($\text{kg}\cdot\text{m}^2\cdot\text{s}^{-1}$)\\
$\mathcal{K}$ & --$\quad$ Constante de separación de Carter ($\text{J}^2\cdot\text{m}^2\cdot\text{s}^{-2}$)\\
$\mathcal{R}(r)$ & --$\quad$ Potencial radial efectivo de Hamilton-Jacobi ($\text{J}^2\cdot\text{m}^2$)\\
$\Phi \equiv \frac{c L_z}{E}$ & --$\quad$ Parámetro de impacto azimutal reducido ($\text{m}$)\\
$\mathcal{Q} \equiv \frac{c^2 \mathcal{K}}{E^2}$ & --$\quad$ Parámetro de impacto polar reducido de Carter ($\text{m}^2$)\\
$b_c \equiv \sqrt{\mathcal{Q}_c + \Phi_c^2}$ & --$\quad$ Parámetro de impacto crítico ($\text{m}$)\\
$r_{\mathrm{ph}}$ & --$\quad$ Radio de la órbita fotónica esférica inestable ($\text{m}$)\\
$R_{\mathrm{eff}}(\delta)$ & --$\quad$ Radio efectivo equivalente de la sombra ($\text{m}$ o $\mu\text{as}$)\\
$D = 1/27$ & --$\quad$ Coeficiente perturbativo de segundo orden del área de la sombra\\
$\gamma_\phi$ & --$\quad$ Exponente orbital de Lyapunov\\
$f_{\mathrm{thin}} \equiv e^{-2\pi\gamma_\phi}$ & --$\quad$ Factor de adelgazamiento exponencial de Sneppen\\
\end{tabular}

\vspace{0.5cm}
\noindent\textbf{Interferometría de muy larga base (VLBI) y observables EHT}
\vspace{0.2cm}

\begin{tabular}{ll}
$(\alpha_c, \beta_c)$ & --$\quad$ Coordenadas celestes en el plano del observador ($\mu\text{as}$)\\
$(u, v)$ & --$\quad$ Coordenadas de frecuencia espacial en el plano de Fourier ($\text{G}\lambda$)\\
$u \equiv \sqrt{u^2 + v^2}$ & --$\quad$ Distancia escalar de la línea de base ($\text{G}\lambda$)\\
$V(u, v)$ & --$\quad$ Visibilidad compleja interferométrica ($\text{Jy}$)\\
$I(\alpha_c, \beta_c)$ & --$\quad$ Distribución de brillo específico en el cielo ($\text{W}\cdot\text{m}^{-2}\cdot\text{Hz}^{-1}\cdot\text{sr}^{-1}$)\\
$F_0$ & --$\quad$ Flujo de emisión total no polarizado ($\text{Jy}$)\\
$w$ & --$\quad$ Ancho FWHM del perfil de emisión del anillo ($\mu\text{as}$)\\
$J_0(x)$ & --$\quad$ Función de Bessel de primer tipo y orden cero\\
$u_1$ & --$\quad$ Frecuencia espacial del primer nulo de visibilidad ($\text{G}\lambda$)\\
$\chi^2_\nu$ & --$\quad$ Chi-cuadrado reducido para la evaluación de bondad de ajuste\\
$\Delta\chi^2_\nu$ & --$\quad$ Diferencia de bondad de ajuste respecto al modelo de Schwarzschild\\
\end{tabular} %% OPCIONAL
%%%%%%%%%%%%%%%%%%%%%%%%%%%%%%%%%%%%%%%%
%       LISTA DE SIGLAS Y ACRÓNIMOS
%%%%%%%%%%%%%%%%%%%%%%%%%%%%%%%%%%%%%%%%
\chapter*{\hfill{\centering Lista de siglas y acrónimos}\hfill}
\addcontentsline{toc}{chapter}{Lista de siglas y acrónimos}
\vspace{0.5cm}
\normalsize

\begin{tabular}{p{2.2cm}l}
\textbf{ALMA}   & --$\quad$ Atacama Large Millimeter/submillimeter Array\\
\textbf{APEX}   & --$\quad$ Atacama Pathfinder Experiment\\
\textbf{BHEX}   & --$\quad$ Black Hole Explorer\\
\textbf{CLEAN}  & --$\quad$ Algoritmo de deconvolución en interferometría de imágenes\\
\textbf{EHT}    & --$\quad$ Event Horizon Telescope (\textit{Telescopio del Horizonte de Sucesos})\\
\textbf{FCF}    & --$\quad$ Facultad de Ciencias Físicas (UNMSM)\\
\textbf{FWHM}   & --$\quad$ Full Width at Half Maximum (\textit{Ancho Completo en la Mitad del Máximo})\\
\textbf{IRAM}   & --$\quad$ Institut de Radioastronomie Millimétrique\\
\textbf{JCMT}   & --$\quad$ James Clerk Maxwell Telescope\\
\textbf{LMT}    & --$\quad$ Large Millimeter Telescope (\textit{Gran Telescopio Milimétrico Alfonso Serrano})\\
\textbf{MAD}    & --$\quad$ Magnetically Arrested Disk (\textit{Disco Magnéticamente Arrestado})\\
\textbf{M87*}   & --$\quad$ Agujero negro supermasivo en la galaxia Messier 87\\
\textbf{ngEHT}  & --$\quad$ next-generation Event Horizon Telescope\\
\textbf{QED}    & --$\quad$ Quantum Electrodynamics (\textit{Electrodinámica Cuántica})\\
\textbf{RG}     & --$\quad$ Relatividad General\\
\textbf{RK4}    & --$\quad$ Algoritmo de integración numérica Runge-Kutta de cuarto orden\\
\textbf{RML}    & --$\quad$ Regularized Maximum Likelihood (\textit{Máxima Verosimilitud Regularizada})\\
\textbf{SANE}   & --$\quad$ Standard and Normal Evolution (\textit{Evolución Estándar y Normal})\\
\textbf{Sgr A*} & --$\quad$ Sagittarius A* (agujero negro supermasivo del Centro Galáctico)\\
\textbf{SI}     & --$\quad$ Sistema Internacional de Unidades\\
\textbf{SMA}    & --$\quad$ Submillimeter Array\\
\textbf{SPT}    & --$\quad$ South Pole Telescope\\
\textbf{UNMSM}  & --$\quad$ Universidad Nacional Mayor de San Marcos\\
\textbf{VLBI}   & --$\quad$ Very Long Baseline Interferometry (\textit{Interferometría de Muy Larga Base})\\
\end{tabular} %% OPCIONAL

\cleardoublepage % Asegura que el Capítulo 1 comience en página impar

\pagenumbering{arabic}
\justifying

%================================================
% AQUÍ INICIA LOS CAPÍTULOS DE LA TESIS
% Adicione los capítulos que crea conveniente 
%================================================

%=============================
%  PRIMER CAPITULO
%=============================
\chapter{INTRODUCCIÓN}

\section{Contexto de la astronomía de precisión: El Event Horizon Telescope (EHT)}
En las últimas décadas, la astrofísica y la relatividad general han ingresado a una era de precisión observacional sin precedentes \citep{EHT_Paper1, EHT_Paper4}. La colaboración internacional del Event Horizon Telescope (EHT) ha logrado capturar las primeras imágenes directas de las sombras y los anillos de emisión alrededor de los agujeros negros supermasivos M87* y Sgr A* \citep{EHT_Paper1, EHT_Paper6}. Estas observaciones abren la oportunidad de testear soluciones de las ecuaciones de Einstein en el régimen de campos gravitacionales fuertes.

% ==============================================================================
% SECCIÓN: MOTIVACIÓN FÍSICA DE LA CARGA ELÉCTRICA EN AGUJEROS NEGROS ASTROFÍSICOS
% ==============================================================================


\section{Motivación física de la carga eléctrica en agujeros negros astrofísicos}

Aunque habitualmente se asume la hipótesis de neutralidad eléctrica estricta ($Q = 0$) para objetos compactos astrofísicos bajo la premisa de una rápida neutralización por el plasma circundante, análisis electrodinámicos y perturbativos en espaciotiempos curvos demuestran la estabilidad y vigencia del estudio de cargas residuales pequeñas $Q \neq 0$ \citep{PhysRevLett.114.151101, s10052-021-09180-y}. En entornos astrofísicos reales, la presencia de campos magnéticos intensos en discos de acreción magnetizados —como los estados de disco magnéticamente arrestado (MAD) modelados para M87* y Sgr A*— induce campos eléctricos por efecto del arrastre de marcos (\textit{frame-dragging}), permitiendo mantener configuraciones de equilibrio con carga residual \citep{EHT_Paper1, 2406.09498v1} Asimismo, procesos de acreción en plasmas ionizados y límites electrodinámicos de descarga por producción de pares restringen cualquier carga remanente estable al régimen de \textbf{carga eléctrica débil}, expresado en unidades del Sistema Internacional (SI) mediante la condición adimensional \citep{1403.5234v2, Broderick_2022_ApJ_927_6}:
\begin{equation}
    \delta \equiv \frac{Q^2}{4\pi\varepsilon_0 G M^2} \ll 1
\end{equation}
En este régimen perturbativo ($\delta \ll 1$), el espaciotiempo de Kerr-Newman constituye un marco teórico y una métrica \textit{proxy} idónea para modelar pequeñas desviaciones geométricas respecto al caso puro de Kerr y contrastarlas con las observaciones de la sombra y el anillo de fotones del \textit{Event Horizon Telescope} \citep{1403.5234v2, Broderick_2022_ApJ_927_6}.

% ==============================================================================
% SECCIÓN: JUSTIFICACIÓN Y OBJETIVOS DE LA INVESTIGACIÓN
% ==============================================================================

\section{Justificación y Objetivos de la Investigación}

El estudio de la sombra de agujeros negros y la estructura fina del anillo de fotones ha cobrado una relevancia observacional sin precedentes tras las publicaciones de la colaboración \textit{Event Horizon Telescope} (EHT) para M87* y Sgr A*. Tradicionalmente, la interpretación de los perfiles de visibilidad interferométrica $|V(u)|$ asume como fondo un espaciotiempo estático de Schwarzschild, atribuyendo cualquier desviación geométrica a incertidumbres en la masa y la distancia ($M/D$) o a la compleja hidrodinámica del plasma de acreción. 

Esta investigación se justifica formalmente al proponer un cambio de paradigma analítico: utilizar como fondo físico exacto el espaciotiempo de Kerr-Newman en el régimen de carga eléctrica débil ($\delta \equiv \frac{Q^2}{4\pi\varepsilon_0 G M^2} \ll 1$), formulado rigurosamente en unidades del Sistema Internacional (SI) y firma métrica de Minkowski $(+,-,-,-)$. Este enfoque permite conectar directamente la electrodinámica en campos fuertes con observables interferométricos del EHT, obteniendo expansiones perturbativas analíticas de segundo orden que aíslan la firma de electrovacío de la carga ($D = 1/27$) frente a modelos de gravedad modificada o agujeros negros regulares. Asimismo, permite contrastar el modelo con los microdatos interferométricos reales ($N = 9\,304$ visibilidades en Sgr A* y $N = 4\,237$ en M87*), proporcionando una interpretación física y epistemológica sobre el rol de $\delta$ como un parámetro de escala efectiva dentro de modelos radiales simétricos.

\subsection{Objetivo General}
Demostrar analíticamente que la modelación de la geodésica de fotones críticos utilizando como espaciotiempo de fondo la métrica de Kerr-Newman bajo la aproximación de carga débil ($\delta \ll 1$), en lugar de una expansión perturbativa simétrica sobre el fondo estático de Schwarzschild, permite derivar correcciones analíticas de segundo orden al radio efectivo de la sombra $R_{\mathrm{eff}}(\delta)$ y minimizar la función de costo estadístico $\Delta\chi^2$ en los perfiles de visibilidad interferométrica $|V(u)|$ medidos por el Event Horizon Telescope (EHT) para Sgr A*, desplazando el primer nulo de visibilidad en un.

\subsection{Objetivos Específicos}
\begin{enumerate}
    \item Deducir el potencial radial efectivo $\mathcal{R}(r)$ y las ecuaciones geodésicas nulas a partir del formalismo de Hamilton-Jacobi y la función de acción $S$ en el espaciotiempo de Kerr-Newman, expresados en el Sistema Internacional (SI) y firma métrica $(+,-,-,-)$, identificando las constantes de movimiento conservadas ($E, L_z, \mathcal{K}$).
    \item Formular las condiciones de órbita fotónica esférica inestable $\mathcal{R}(r_{\mathrm{ph}}) = 0$ y $\left.\frac{d\mathcal{R}}{dr}\right|_{r_{\mathrm{ph}}} = 0$ en función de los parámetros de impacto reducidos $\Phi \equiv \frac{c L_z}{E}$ y $\mathcal{Q} \equiv \frac{c^2 \mathcal{K}}{E^2}$, derivando las funciones críticas $(\Phi_c, \mathcal{Q}_c)$ que determinan la sombra celeste $(\alpha_c, \beta_c)$ y obteniendo las expansiones de Taylor analíticas para el radio fotónico $r_{\mathrm{ph}}(\delta)$, el radio efectivo $R_{\mathrm{eff}}(\delta)$ (con coeficiente de segundo orden $-7/216$) y la posición del primer nulo de visibilidad $u_1(\delta)$ (con coeficiente $+13/216$).
    \item Analizar la ecuación diferencial del oscilador inestable de Lyapunov y demostrar cómo la degradación del exponente orbital por efecto del momento angular ($\gamma_\phi < 1$) relaja el factor de adelgazamiento de Sneppen ($f_{\mathrm{thin}} = e^{-2\pi\gamma_\phi}$), justificando el ensanchamiento orbital del primer subanillo de fotones $n=1$.
    \item Adaptar el formalismo perturbativo de Kobialko para la expansión del parámetro de impacto crítico $b_c^2 = \mathcal{Q}_c + \Phi_c^2$ a métricas rotacionales no diagonales, aislando analíticamente la firma unívoca de carga electrostática $D = 1/27$ frente a soluciones de agujeros negros regulares (Bardeen $2/27$ y Hayward $8/243$).
    \item Ejecutar un protocolo de ajuste de mínimos cuadrados sobre el modelo de visibilidad de Bessel $J_0$ con envolvente gaussiana empleando las visibilidades reales del EHT para Sgr A* ($2024\text{-D02-01}$) y M87* ($2024\text{-D01-01}$), contrastando la variación en la bondad de ajuste ($\Delta\chi^2_\nu$) frente a las incertidumbres sistemáticas en la razón masa-distancia ($M/D$).
\end{enumerate} %% 1o capítulo
%=============================
%  SEGUNDO CAPITULO
%=============================
\chapter{FUNDAMENTOS Y MÉTRICAS DE AGUJEROS NEGROS EN EL SISTEMA INTERNACIONAL}

\section{Ecuaciones de Einstein-Maxwell }

En el Sistema Internacional de Unidades (SI), la dinámica acoplada entre la geometría espaciotemporal tetradimensional $\mathcal{M}$ y el campo electromagnético está gobernada por la Acción de Einstein-Maxwell \citep{M_P_Hobson_General_Relativity, Sean_Carroll_Spacetime, 05villegas}:
\begin{equation}
    S[g^{\mu\nu}, A_\alpha] = \int_{\mathcal{M}} \left( \frac{c^4}{16\pi G} R - \frac{1}{4\mu_0} F_{\alpha\beta}F^{\alpha\beta} \right) \sqrt{-g} \, d^4x + S_m
    \label{eq:accion_EM_SI}
\end{equation}
donde $G$ es la constante de gravitación universal, $c$ es la velocidad de la luz en el vacío, $\mu_0$ es la permeabilidad magnética del vacío, $g \equiv \det(g_{\mu\nu}) < 0$ es el determinante de la métrica, $R = g^{\mu\nu} R_{\mu\nu}$ es el escalar de Ricci, y $F_{\alpha\beta} \equiv \partial_\alpha A_\beta - \partial_\beta A_\alpha$ es el tensor electromagnético expresado en términos del cuadripotencial $A_\alpha$. El elemento de línea plano en coordenadas esféricas adopta la forma $ds^2 = c^2 dt^2 - dr^2 - r^2 d\theta^2 - r^2 \sin^2\theta d\phi^2$ \citep{M_P_Hobson_General_Relativity, 05villegas}.

Para obtener las ecuaciones de campo de Einstein-Maxwell, aplicamos el principio de mínima acción variando $S$ respecto a la métrica contravariante $g^{\mu\nu}$, manteniendo fijo el campo electromagnético $A_\alpha$. En primer lugar, utilizando la identidad de Jacobi $\delta g = g \, g^{\mu\nu} \delta g_{\mu\nu} = -g \, g_{\mu\nu} \delta g^{\mu\nu}$, la variación del elemento de volumen $\sqrt{-g}$ resulta:
\begin{equation}
    \delta \sqrt{-g} = -\frac{1}{2\sqrt{-g}} \delta g
    \label{eq:var_vol_1}
\end{equation}
\begin{equation}
    \delta \sqrt{-g} = -\frac{1}{2} \sqrt{-g} \, g_{\mu\nu} \delta g^{\mu\nu}
    \label{eq:var_vol_2}
\end{equation}

A continuación, calculamos la variación de la parte gravitacional de la acción, $S_{\mathrm{EH}} = \frac{c^4}{16\pi G} \int R \sqrt{-g} \, d^4x$:
\begin{equation}
    \delta S_{\mathrm{EH}} = \frac{c^4}{16\pi G} \int \left( R_{\mu\nu} \delta g^{\mu\nu} \sqrt{-g} + g^{\mu\nu} \delta R_{\mu\nu} \sqrt{-g} + R \, \delta \sqrt{-g} \right) d^4x
    \label{eq:var_seh_1}
\end{equation}
Por la identidad de Palatini, el término $g^{\mu\nu} \delta R_{\mu\nu} = \nabla_\alpha (\nabla_\beta \delta g^{\alpha\beta} - g^{\alpha\beta} \nabla_\gamma \delta g^{\gamma}_{\phantom{\gamma}\alpha})$ constituye una divergencia total que se anula en la frontera del espacio-tiempo mediante el teorema de Gauss. Sustituyendo la ec.~(\ref{eq:var_vol_2}) en la ec.~(\ref{eq:var_seh_1}), la variación del término de Einstein-Hilbert toma la forma:
\begin{equation}
    \delta S_{\mathrm{EH}} = \frac{c^4}{16\pi G} \int \left( R_{\mu\nu} - \frac{1}{2} g_{\mu\nu} R \right) \delta g^{\mu\nu} \sqrt{-g} \, d^4x
    \label{eq:var_seh_2}
\end{equation}

Por otro lado, para variar el término electromagnético expresamos la cantidad escalar como $F_{\alpha\beta} F^{\alpha\beta} = g^{\alpha\gamma} g^{\beta\delta} F_{\alpha\beta} F_{\gamma\delta}$. La variación directa respecto a la métrica $g^{\mu\nu}$ da:
\begin{equation}
    \delta \left( F_{\alpha\beta} F^{\alpha\beta} \right) = \delta(g^{\alpha\gamma}) g^{\beta\delta} F_{\alpha\beta} F_{\gamma\delta} + g^{\alpha\gamma} \delta(g^{\beta\delta}) F_{\alpha\beta} F_{\gamma\delta}
    \label{eq:var_f2_1}
\end{equation}
\begin{equation}
    \delta \left( F_{\alpha\beta} F^{\alpha\beta} \right) = 2 F_{\mu\alpha} F_{\nu}^{\phantom{\nu}\alpha} \, \delta g^{\mu\nu}
    \label{eq:var_f2_2}
\end{equation}
Con este resultado, la variación de la acción electromagnética $S_{\mathrm{EM}} = -\frac{1}{4\mu_0} \int F_{\alpha\beta} F^{\alpha\beta} \sqrt{-g} \, d^4x$ se escribe como:
\begin{equation}
    \delta S_{\mathrm{EM}} = -\frac{1}{4\mu_0} \int \left[ \delta\left(F_{\alpha\beta} F^{\alpha\beta}\right) \sqrt{-g} + F_{\alpha\beta} F^{\alpha\beta} \delta\sqrt{-g} \right] d^4x
    \label{eq:var_sem_1}
\end{equation}
Sustituyendo las ecs.~(\ref{eq:var_vol_2}) y~(\ref{eq:var_f2_2}) en la ec.~(\ref{eq:var_sem_1}):
\begin{equation}
    \delta S_{\mathrm{EM}} = -\frac{1}{4\mu_0} \int \left[ 2 F_{\mu\alpha} F_{\nu}^{\phantom{\nu}\alpha} - \frac{1}{2} g_{\mu\nu} F_{\alpha\beta} F^{\alpha\beta} \right] \delta g^{\mu\nu} \sqrt{-g} \, d^4x
    \label{eq:var_sem_2}
\end{equation}
\begin{equation}
    \delta S_{\mathrm{EM}} = -\frac{1}{2} \int \left[ \frac{1}{\mu_0} \left( F_{\mu\,\alpha}F_\nu^{\phantom{\nu}\alpha} - \frac{1}{4}g_{\mu\nu} F_{\alpha\beta}F^{\alpha\beta} \right) \right] \delta g^{\mu\nu} \sqrt{-g} \, d^4x
    \label{eq:var_sem_3}
\end{equation}
Esta última expresión nos permite definir el tensor de energía-impulso electromagnético en el SI como \citep{M_P_Hobson_General_Relativity, 05villegas}:
\begin{equation}
    T_{\mu\nu}^{(\mathrm{EM})} \equiv -\frac{2}{\sqrt{-g}} \frac{\delta S_{\mathrm{EM}}}{\delta g^{\mu\nu}}
    \label{eq:def_tensor_EM}
\end{equation}
\begin{equation}
    T_{\mu\nu}^{(\mathrm{EM})} = \frac{1}{\mu_0} \left( F_{\mu\,\alpha}F_\nu^{\phantom{\nu}\alpha} - \frac{1}{4}g_{\mu\nu} F_{\alpha\beta}F^{\alpha\beta} \right)
    \label{eq:tensor_EM_SI}
\end{equation}

Exigiendo la condición de punto estacionario $\delta S = \delta S_{\mathrm{EH}} + \delta S_{\mathrm{EM}} = 0$ para cualquier variación arbitraria $\delta g^{\mu\nu}$, obtenemos las ecuaciones de campo de Einstein en el SI \citep{M_P_Hobson_General_Relativity, Sean_Carroll_Spacetime}:
\begin{equation}
    R_{\mu\nu} - \frac{1}{2}g_{\mu\nu}R = \frac{8\pi G}{c^4} T_{\mu\nu}^{(\mathrm{EM})}
    \label{eq:einstein_SI}
\end{equation}
Para simplificar esta ecuación, calculamos la traza del tensor de energía-impulso $T^{(\mathrm{EM})} \equiv g^{\mu\nu} T_{\mu\nu}^{(\mathrm{EM})}$ en cuatro dimensiones:
\begin{equation}
    T^{(\mathrm{EM})} = \frac{1}{\mu_0} \left( g^{\mu\nu} F_{\mu\alpha} F_{\nu}^{\phantom{\nu}\alpha} - \frac{1}{4} g^{\mu\nu} g_{\mu\nu} F_{\alpha\beta} F^{\alpha\beta} \right)
    \label{eq:traza_tem_1}
\end{equation}
Dado que $g^{\mu\nu} g_{\mu\nu} = 4$, la ec.~(\ref{eq:traza_tem_1}) se reduce a:
\begin{equation}
    T^{(\mathrm{EM})} = \frac{1}{\mu_0} \left( F_{\alpha\beta} F^{\alpha\beta} - \frac{4}{4} F_{\alpha\beta} F^{\alpha\beta} \right) \equiv 0
    \label{eq:traza_tem_2}
\end{equation}
Al ser $T^{(\mathrm{EM})} = 0$, contratando la ec.~(\ref{eq:einstein_SI}) con $g^{\mu\nu}$ obtenemos $R - \frac{1}{2}(4)R = 0 \implies -R = 0 \implies R = 0$. Por lo tanto, en el electrovacío el escalar de Ricci se anula idénticamente y las ecuaciones de Einstein adoptan la forma simplificada:
\begin{equation}
    R_{\mu\nu} = \frac{8\pi G}{c^4} T_{\mu\nu}^{(\mathrm{EM})}
    \label{eq:einstein_electrovacio}
\end{equation}

Finalmente, derivamos las ecuaciones de Maxwell variando la acción $S_{\mathrm{EM}}$ respecto al cuadripotencial $A_\beta$. Utilizando la relación $\delta F_{\alpha\beta} = \partial_\alpha \delta A_\beta - \partial_\beta \delta A_\alpha = 2 \nabla_{[\alpha} \delta A_{\beta]}$:
\begin{equation}
    \delta_A S_{\mathrm{EM}} = -\frac{1}{\mu_0} \int F^{\alpha\beta} (\nabla_\alpha \delta A_\beta) \sqrt{-g} \, d^4x
    \label{eq:var_maxwell_1}
\end{equation}
Integrando por partes y aplicando la compatibilidad métrica $\nabla_\alpha \sqrt{-g} = 0$:
\begin{equation}
    \delta_A S_{\mathrm{EM}} = \frac{1}{\mu_0} \int \left( \nabla_\alpha F^{\alpha\beta} \right) \delta A_\beta \sqrt{-g} \, d^4x = 0
    \label{eq:var_maxwell_2}
\end{equation}
De la arbitrariedad de $\delta A_\beta$ se derivan las ecuaciones de Maxwell covariantes en el vacío \citep{M_P_Hobson_General_Relativity, Sean_Carroll_Spacetime}:
\begin{equation}
    \nabla_\alpha F^{\alpha\beta} = 0
    \label{eq:maxwell_covariante}
\end{equation}
que en términos de derivadas parciales se expresan como:
\begin{equation}
    \frac{1}{\sqrt{-g}} \partial_\alpha \left( \sqrt{-g} \, F^{\alpha\beta} \right) = 0
    \label{eq:maxwell_parciales}
\end{equation}
junto con la identidad de Bianchi electromagnética:
\begin{equation}
    \partial_\alpha F_{\beta\gamma} + \partial_\beta F_{\gamma\alpha} + \partial_\gamma F_{\alpha\beta} = 0
    \label{eq:bianchi_em}
\end{equation}
El conjunto acoplado de las ecuaciones~(\ref{eq:einstein_electrovacio}), (\ref{eq:maxwell_parciales}) y~(\ref{eq:bianchi_em}) constituye el sistema de Einstein-Maxwell  cuya solución exacta e incompresible para un agujero negro con masa $M$, momento angular $J$ y carga $Q$ es la métrica de Kerr-Newman \citep{1403.5234v2, 05villegas}.




\section{La métrica de Kerr-Newman en unidades SI}

Al incorporar el momento angular $J$ del objeto compacto, la rotación destruye la simetría esférica pero conserva la simetría axial respecto al eje de rotación y la estacionariedad temporal \citep{1403.5234v2, Sean_Carroll_Spacetime}. El parámetro de espín con dimensiones de longitud se define en el SI como \citep{M_P_Hobson_General_Relativity, s10052-021-09180-y}:
\begin{equation}
    a \equiv \frac{J}{M c}
    \label{eq:parametro_espin_SI}
\end{equation}

En coordenadas de Boyer-Lindquist $(t, r, \theta, \phi)$ y bajo la convención de firma de Minkowski $(+,-,-,-)$, el elemento de línea exacto del espaciotiempo de Kerr-Newman se expresa en el SI como \citep{1403.5234v2, pdfcoffee.com_chandrasekhar}:
\begin{equation}
    ds^2 = g_{tt} c^2 dt^2 + 2 c \, g_{t\phi} dt d\phi + g_{rr} dr^2 + g_{\theta\theta} d\theta^2 + g_{\phi\phi} d\phi^2
    \label{eq:linea_kerr_newman_general}
\end{equation}
donde las componentes covariantes del tensor métrico $g_{\mu\nu}$ son \citep{1403.5234v2, s10052-021-09180-y}:
\begin{equation}
    g_{tt}(r,\theta) = 1 - \frac{2r_g r - r_g^2 \delta}{\Sigma(r,\theta)}
    \label{eq:gtt_KN}
\end{equation}
\begin{equation}
    g_{t\phi}(r,\theta) = \frac{a \sin^2\theta}{\Sigma(r,\theta)} \left( 2r_g r - r_g^2 \delta \right)
    \label{eq:gtphi_KN}
\end{equation}
\begin{equation}
    g_{rr}(r,\theta) = -\frac{\Sigma(r,\theta)}{\Delta(r)}
    \label{eq:grr_KN}
\end{equation}
\begin{equation}
    g_{\theta\theta}(r,\theta) = -\Sigma(r,\theta)
    \label{eq:gthetatheta_KN}
\end{equation}
\begin{equation}
    g_{\phi\phi}(r,\theta) = -\left[ r^2 + a^2 + \frac{a^2 \sin^2\theta}{\Sigma(r,\theta)} \left( 2r_g r - r_g^2 \delta \right) \right] \sin^2\theta
    \label{eq:gphiphi_KN}
\end{equation}
Las funciones métricas auxiliares $\Sigma(r,\theta)$ y $\Delta(r)$ se definen explícitamente como \citep{1403.5234v2, s10052-021-09180-y}:
\begin{equation}
    \Sigma(r,\theta) = r^2 + a^2 \cos^2\theta
    \label{eq:sigma_def_detallada}
\end{equation}
\begin{equation}
    \Delta(r) = r^2 - \frac{2GM}{c^2}r + a^2 + \frac{G Q^2}{4\pi\varepsilon_0 c^4} = r^2 - 2 r_g r + a^2 + r_g^2 \delta
    \label{eq:delta_def_detallada}
\end{equation}

El cuadripotencial electromagnético acoplado a la geometría de Kerr-Newman viene dado en coordenadas de Boyer-Lindquist por \citep{1403.5234v2, Sean_Carroll_Spacetime}:
\begin{equation}
    A_\mu = \frac{Q r}{4\pi\varepsilon_0 \Sigma c} \left( 1, 0, 0, -a \sin^2\theta \right)
    \label{eq:cuadripotencial_KN}
\end{equation}
donde la componente temporal $A_t = \frac{Q r}{4\pi\varepsilon_0 \Sigma c}$ representa el potencial eléctrico modificado por la rotación y $A_\phi = -\frac{Q r a \sin^2\theta}{4\pi\varepsilon_0 \Sigma c}$ describe el campo magnético dipolar inducido por el arrastre del momento angular \citep{M_P_Hobson_General_Relativity}.

El determinante de la métrica de Kerr-Newman se calcula directamente a partir del bloque no diagonal $(t,\phi)$ y los componentes diagonales $(r,\theta)$:
\begin{equation}
    g \equiv \det(g_{\mu\nu}) = (g_{tt} g_{\phi\phi} - g_{t\phi}^2) g_{rr} g_{\theta\theta} = -\Sigma^2 \sin^2\theta
    \label{eq:det_kerr_newman}
\end{equation}
\begin{equation}
    \sqrt{-g} = \Sigma \sin^2\theta \dots \implies \sqrt{-g} = \Sigma \sin\theta
    \label{eq:raiz_det_kerr_newman}
\end{equation}

Invertimos la matriz métrica covariante $g_{\mu\nu}$ para obtener los componentes contravariantes $g^{\mu\nu}$, fundamentales para la formulación del potencial radial efectivo mediante la ecuación de Hamilton-Jacobi en el Capítulo III \citep{pdfcoffee.com_chandrasekhar, 1403.5234v2}:
\begin{equation}
    g^{tt}(r,\theta) = \frac{(r^2 + a^2)^2 - a^2 \Delta(r) \sin^2\theta}{\Sigma(r,\theta) \, \Delta(r)}
    \label{eq:g_contra_tt}
\end{equation}
\begin{equation}
    g^{t\phi}(r,\theta) = \frac{a \left[ (r^2 + a^2) - \Delta(r) \right]}{\Sigma(r,\theta) \, \Delta(r)} = \frac{a \left( 2 r_g r - r_g^2 \delta \right)}{\Sigma(r,\theta) \, \Delta(r)}
    \label{eq:g_contra_tphi}
\end{equation}
\begin{equation}
    g^{rr}(r,\theta) = -\frac{\Delta(r)}{\Sigma(r,\theta)}
    \label{eq:g_contra_rr}
\end{equation}
\begin{equation}
    g^{\theta\theta}(r,\theta) = -\frac{1}{\Sigma(r,\theta)}
    \label{eq:g_contra_thetatheta}
\end{equation}
\begin{equation}
    g^{\phi\phi}(r,\theta) = -\frac{\Delta(r) - a^2 \sin^2\theta}{\Sigma(r,\theta) \, \Delta(r) \sin^2\theta}
    \label{eq:g_contra_phiphi}
\end{equation}

La estructura global de la métrica de Kerr-Newman posee dos superficies características fundamentales:
\begin{enumerate}
    \item \textbf{Horizontes de eventos ($r_\pm$):} Corresponden a las hipersuperficies nulas dadas por las raíces de la función $\Delta(r) = 0$ \citep{1403.5234v2, s10052-021-09180-y}:
    \begin{equation}
        r_\pm = r_g \left( 1 \pm \sqrt{1 - \frac{a^2}{r_g^2} - \delta} \right)
        \label{eq:horizontes_KN}
    \end{equation}
    \item \textbf{Ergosfera o Límite Estático ($r_{\text{ergo}}(\theta)$):} Corresponde al lugar geométrico donde el componente $g_{tt}(r,\theta) = 0$, lo que impide la existencia de observadores estáticos con cuadrivelocidad paralela al vector de Killing temporal \citep{M_P_Hobson_General_Relativity, Sean_Carroll_Spacetime}:
    \begin{equation}
        r_{\text{ergo}}(\theta) = r_g \left( 1 + \sqrt{1 - \frac{a^2 \cos^2\theta}{r_g^2} - \delta} \right)
        \label{eq:ergosfera_KN}
    \end{equation}
\end{enumerate}
La condición para la existencia de un agujero negro no extremall exige que el discriminante sea estrictamente positivo: $a^2/r_g^2 + \delta < 1$. En el régimen de carga débil ($\delta \ll 1$), la métrica de Kerr-Newman~(\ref{eq:linea_kerr_newman_general}) constituye la perturbación exacta adecuada para analizar el impacto de la carga electrostática sobre la dinámica de fotones sin alterar la integrabilidad matemática del espaciotiempo \citep{1403.5234v2, s10052-021-09180-y}. %% 2o capítulo
\chapter{DINÁMICA GEODÉSICA Y PARÁMETRO DE IMPACTO EN KERR-NEWMAN}

% ==============================================================================
% SECCIÓN: DEDUCCIÓN DEL LAGRANGIANO Y CONSTANTES DE MOVIMIENTO
% ==============================================================================

\section{Deducción del Lagrangiano y constantes de movimiento}

Siguiendo el formalismo analítico de la mecánica clásica y la relatividad general \citep{05villegas, M_P_Hobson_General_Relativity, Sean_Carroll_Spacetime}, la dinámica geodésica de partículas de prueba y fotones en el espaciotiempo de Kerr-Newman se deriva de forma compacta a través de la densidad lagrangiana por unidad de masa:
\begin{equation}
    \mathcal{L}(x^\mu, \dot{x}^\mu) = \frac{1}{2} g_{\mu\nu}(x) \frac{dx^\mu}{d\lambda} \frac{dx^\nu}{d\lambda} = \frac{1}{2} g_{\mu\nu}(x) \dot{x}^\mu \dot{x}^\nu
    \label{eq:lagrangiano_gen}
\end{equation}
donde $\lambda$ representa un parámetro afín adecuado a lo largo de la trayectoria (el tiempo propio $\tau$ para partículas masivas o un parámetro afín arbitrario para fotones con $m=0$), y el punto denota la derivada respecto a $\lambda$ ($\dot{x}^\mu \equiv dx^\mu/d\lambda$).

\subsection{Simetrías del espaciotiempo y teorema de Noether}

La métrica de Kerr-Newman en coordenadas de Boyer-Lindquist $(t, r, \theta, \phi)$ no depende explícitamente ni del tiempo $t$ ni de la coordenada azimutal $\phi$ \citep{1403.5234v2, pdfcoffee.com_chandrasekhar}. Por consiguiente, la geometría admite dos vectores de Killing independientes:
\begin{enumerate}
    \item \textbf{Vector de Killing de estacionariedad ($\xi^\mu$):} $\xi^\mu = (1, 0, 0, 0)$, asociado a la invarianza por traducciones temporales ($\partial_t g_{\mu\nu} = 0$).
    \item \textbf{Vector de Killing de axisimetría ($\eta^\mu$):} $\eta^\mu = (0, 0, 0, 1)$, asociado a la invarianza por rotaciones alrededor del eje de simetría ($\partial_\phi g_{\mu\nu} = 0$).
\end{enumerate}

Por el teorema de Noether, las coordenadas $t$ y $\phi$ son coordenadas cíclicas en el Lagrangiano $\mathcal{L}$, lo que garantiza la existencia de dos impulsos canónicamente conjugados que se conservan idénticamente a lo largo de toda geodésica \citep{05villegas, M_P_Hobson_General_Relativity}.

\subsection{Matriz métrica en el plano ecuatorial \texorpdfstring{($\theta = \pi/2$)}{}}

Para el estudio de las órbitas fotónicas en el plano ecuatorial ($\theta = \pi/2$), imponemos las condiciones geométricas $\dot{\theta} = 0$, $\sin\theta = 1$ y $\cos\theta = 0$. Bajo esta restricción, la función auxiliar de Boyer-Lindquist se reduce a $\Sigma(r,\pi/2) = r^2$, y la función del horizonte adopta la forma compacta $\Delta(r) = r^2 - 2 r_g r + a^2 + r_g^2 \delta$, donde $r_g \equiv \frac{GM}{c^2}$ es el radio gravitacional y $\delta \equiv \frac{Q^2}{4\pi\varepsilon_0 G M^2} \ll 1$ es el parámetro adimensional de carga eléctrica débil \citep{s10052-021-09180-y, 1403.5234v2}.

La matriz covariante del tensor métrico $(g_{\mu\nu})$ en la base $(c t, r, \theta, \phi)$ para la firma de Minkowski $(+,-,-,-)$ en el plano ecuatorial se escribe explícitamente como:
\begin{equation}
    (g_{\mu\nu}) = \begin{pmatrix}
        1 - \frac{2r_g}{r} + \frac{r_g^2 \delta}{r^2} & 0 & 0 & \frac{a}{r}\left(\frac{2r_g}{r} - \frac{r_g^2 \delta}{r^2}\right) \\[0.2cm]
        0 & -\frac{r^2}{\Delta(r)} & 0 & 0 \\[0.2cm]
        0 & 0 & -r^2 & 0 \\[0.2cm]
        \frac{a}{r}\left(\frac{2r_g}{r} - \frac{r_g^2 \delta}{r^2}\right) & 0 & 0 & -\left[ r^2 + a^2 + \frac{a^2}{r}\left(\frac{2r_g}{r} - \frac{r_g^2 \delta}{r^2}\right) \right]
    \end{pmatrix}
    \label{eq:matriz_g_KN}
\end{equation}

Sustituyendo las componentes de la matriz~(\ref{eq:matriz_g_KN}) en la ec.~(\ref{eq:lagrangiano_gen}), desplegamos la expresión del Lagrangiano $\mathcal{L}$ en el plano ecuatorial:
\begin{align}
    \mathcal{L} &= \frac{1}{2} \left[ g_{tt} c^2 \dot{t}^2 + 2 c \, g_{t\phi} \dot{t} \dot{\phi} + g_{rr} \dot{r}^2 + g_{\phi\phi} \dot{\phi}^2 \right] \nonumber \\[0.2cm]
    &= \frac{1}{2} \left[ \left(1 - \frac{2r_g}{r} + \frac{r_g^2 \delta}{r^2}\right) c^2 \dot{t}^2 + \frac{2 a c}{r}\left(\frac{2r_g}{r} - \frac{r_g^2 \delta}{r^2}\right) \dot{t}\dot{\phi} \right. \nonumber \\
    &\quad \left. - \frac{r^2}{\Delta(r)} \dot{r}^2 - \left(r^2 + a^2 + \frac{a^2}{r}\left(\frac{2r_g}{r} - \frac{r_g^2 \delta}{r^2}\right)\right) \dot{\phi}^2 \right]
    \label{eq:L_desarrollado}
\end{align}

\subsection{Deducción analítica de las constantes de movimiento}

Aplicamos el formalismo de Euler-Lagrange para cada una de las coordenadas cíclicas:

\paragraph{1. Conservación de la energía ($E$):}
Para la coordenada temporal $q^0 = t$, la ecuación de Euler-Lagrange establece:
\begin{equation}
    \frac{d}{d\lambda}\left( \frac{\partial \mathcal{L}}{\partial \dot{t}} \right) - \frac{\partial \mathcal{L}}{\partial t} = 0
    \label{eq:euler_lagrange_t}
\end{equation}
Dado que $\frac{\partial \mathcal{L}}{\partial t} = 0$, se deduce que el momento conjugado temporal $p_t$ es constante a lo largo de la trayectoria. Derivando directamente la ec.~(\ref{eq:L_desarrollado}) respecto a $\dot{t}$:
\begin{equation}
    p_t \equiv \frac{\partial \mathcal{L}}{\partial \dot{t}} = g_{tt} c^2 \dot{t} + g_{t\phi} c \, \dot{\phi} = \text{constante} \equiv E
    \label{eq:E_conservada_paso1}
\end{equation}
\begin{equation}
    E = \left(1 - \frac{2r_g}{r} + \frac{r_g^2 \delta}{r^2}\right) c^2 \dot{t} + \frac{a c}{r}\left(\frac{2r_g}{r} - \frac{r_g^2 \delta}{r^2}\right) \dot{\phi}
    \label{eq:E_conservada_def}
\end{equation}
donde $E$ representa la energía conservada por el fotón tal como la mide un observador en el infinito.

\paragraph{2. Conservación del momento angular azimutal ($L_z$):}
Para la coordenada azimutal $q^3 = \phi$, la ecuación de Euler-Lagrange exige:
\begin{equation}
    \frac{d}{d\lambda}\left( \frac{\partial \mathcal{L}}{\partial \dot{\phi}} \right) - \frac{\partial \mathcal{L}}{\partial \phi} = 0
    \label{eq:euler_lagrange_phi}
\end{equation}
Al ser $\frac{\partial \mathcal{L}}{\partial \phi} = 0$, el momento conjugado azimutal $p_\phi$ se conserva idénticamente. Derivando la ec.~(\ref{eq:L_desarrollado}) respecto a $\dot{\phi}$:
\begin{equation}
    p_\phi \equiv \frac{\partial \mathcal{L}}{\partial \dot{\phi}} = g_{t\phi} c \, \dot{t} + g_{\phi\phi} \dot{\phi} = \text{constante} \equiv -L_z
    \label{eq:Lz_conservada_paso1}
\end{equation}
\begin{equation}
    -L_z = \frac{a c}{r}\left(\frac{2r_g}{r} - \frac{r_g^2 \delta}{r^2}\right) \dot{t} - \left[ r^2 + a^2 + \frac{a^2}{r}\left(\frac{2r_g}{r} - \frac{r_g^2 \delta}{r^2}\right) \right] \dot{\phi}
    \label{eq:Lz_conservada_def}
\end{equation}
El signo menos en el momento angular $p_\phi = -L_z$ proviene de la firma métrica espacial $(-,-,-)$ de la convención de Minkowski $(+,-,-,-)$ \citep{05villegas, M_P_Hobson_General_Relativity}. Las ecuaciones~(\ref{eq:E_conservada_def}) y~(\ref{eq:Lz_conservada_def}) constituyen un sistema lineal cerrado de dos ecuaciones algebraicas para las velocidades coordenadas $(\dot{t}, \dot{\phi})$, cuya resolución analítica permite desacoplar los momentos y formular el potencial radial efectivo $\mathcal{R}(r)$ en el Capítulo III.

\section{Separabilidad mediante el formalismo de Hamilton-Jacobi, la Constante de Carter y el potencial radial \texorpdfstring{$\mathcal{R}(r)$}{R(r)}}

Para describir la dinámica de fotones a partir de sus fundamentos analíticos, recurrimos a la Ecuación de Hamilton-Jacobi \citep{pdfcoffee.com_chandrasekhar, 05villegas}. En un espaciotiempo curvado, el Hamiltoniano $H(x^\mu, p_\mu)$ para una partícula de prueba se construye a partir de las componentes contravariantes del tensor métrico $g^{\mu\nu}$ y de los momentos canónicos conjugados $p_\mu = g_{\mu\nu} \dot{x}^\nu$ mediante la expresión \citep{05villegas}:
\begin{equation}
    H(x^\mu, p_\mu) = \frac{1}{2} g^{\mu\nu} p_\mu p_\nu
    \label{eq:hamiltoniano_gen}
\end{equation}
Para la propagación de fotones (geodésicas nulas), la condición de norma $g_{\mu\nu}\dot{x}^\mu\dot{x}^\nu = 0$ exige que el Hamiltoniano satisfaga la restricción $H = 0$ \citep{05villegas, chaos_fractals_photon_ring}.

Sustituyendo los momentos conjugados $p_\mu$ por las derivadas parciales de la Función Principal de Hamilton o Acción de Jacobi $S(x^\mu)$, mediante la relación $p_\mu = \frac{\partial S}{\partial x^\mu}$, obtenemos la Ecuación de Hamilton-Jacobi \citep{05villegas}:
\begin{equation}
    g^{\mu\nu} \frac{\partial S}{\partial x^\mu} \frac{\partial S}{\partial x^\nu} = 0
    \label{eq:HJ_general}
\end{equation}
\newpage
Para evaluar explícitamente la Ecuación~\eqref{eq:HJ_general}, tomamos las ecuaciones ~\eqref{eq:g_contra_tt},~\eqref{eq:g_contra_tphi},~\eqref{eq:g_contra_rr}, ~\eqref{eq:g_contra_thetatheta}, ~\eqref{eq:g_contra_phiphi} del tensor métrico de Kerr-Newman $g_{\mu\nu}$ en coordenadas de Boyer-Lindquist $(t, r, \theta, \phi)$ $g^{\mu\nu}$ \citep{1403.5234v2, chaos_fractals_photon_ring}:


Desarrollando la Ecuación~\eqref{eq:HJ_general} término a término mediante las componentes contravariantes~(\ref{eq:g_contra_tt}--\ref{eq:g_contra_phiphi}):
\begin{equation}
    g^{tt}\left(\frac{\partial S}{\partial t}\right)^2 + 2g^{t\phi}\left(\frac{\partial S}{\partial t}\right)\left(\frac{\partial S}{\partial \phi}\right) + g^{\phi\phi}\left(\frac{\partial S}{\partial \phi}\right)^2 + g^{rr}\left(\frac{\partial S}{\partial r}\right)^2 + g^{\theta\theta}\left(\frac{\partial S}{\partial \theta}\right)^2 = 0
    \label{eq:HJ_expandida}
\end{equation}

Puesto que las componentes de la métrica de Kerr-Newman no dependen de las coordenadas $t$ ni $\phi$, estas son cíclicas. La simetría por traslación temporal y rotación azimutal garantiza la existencia de las constantes de movimiento $E$ (energía conservada) y $L_z$ (momento angular azimutal) \citep{ chaos_fractals_photon_ring}:
\begin{equation}
    p_t = \frac{\partial S}{\partial t} \equiv -E, \quad p_\phi = \frac{\partial S}{\partial \phi} \equiv L_z
    \label{eq:constantes_E_Lz}
\end{equation}


\subsection{Impulsos canónicos y separación de la ecuación de Hamilton-Jacobi}

En el formalismo de Hamilton-Jacobi para geodésicas nulas ($m=0$), la función de acción separada se propone como \citep{pdfcoffee.com_chandrasekhar, 1403.5234v2}:
\begin{equation}
    S(t,r,\theta,\phi) = -E t + L_z \phi + S_r(r) + S_\theta(\theta)
    \label{eq:ansatz_S_oficial}
\end{equation}
donde las derivadas parciales de la acción definen los impulsos canónicamente conjugados $p_\mu \equiv \frac{\partial S}{\partial x^\mu}$. Empleando la relación entre los impulsos y las cuadrivelocidades $p_\mu = g_{\mu\nu} \dot{x}^\nu$, explicitamos cada componente en el Sistema Internacional (SI) bajo la firma de Minkowski $(+,-,-,-)$:
\begin{enumerate}
    \item \textbf{Impulso temporal:} $p_t \equiv \frac{\partial S}{\partial t} = -E = g_{tt} c^2 \dot{t} + g_{t\phi} c \, \dot{\phi}$
    \item \textbf{Impulso azimutal:} $p_\phi \equiv \frac{\partial S}{\partial \phi} = L_z = g_{t\phi} c \, \dot{t} + g_{\phi\phi} \dot{\phi}$
    \item \textbf{Impulso radial:} $p_r \equiv \frac{\partial S}{\partial r} = \frac{dS_r}{dr} = g_{rr} \dot{r} = -\frac{\Sigma(r,\theta)}{\Delta(r)} \dot{r}$
    \item \textbf{Impulso angular:} $p_\theta \equiv \frac{\partial S}{\partial \theta} = \frac{dS_\theta}{d\theta} = g_{\theta\theta} \dot{\theta} = -\Sigma(r,\theta) \dot{\theta}$
\end{enumerate}

Sustituyendo los impulsos $p_\mu$ en la ecuación covariante de Hamilton-Jacobi $g^{\mu\nu} p_\mu p_\nu = 0$, obtenemos:
\begin{equation}
    g^{tt} E^2 - 2g^{t\phi} E L_z + g^{\phi\phi} L_z^2 + g^{rr} \left(\frac{dS_r}{dr}\right)^2 + g^{\theta\theta} \left(\frac{dS_\theta}{d\theta}\right)^2 = 0
    \label{eq:HJ_impulsos}
\end{equation}
Sustituyendo las componentes métricas espaciales contravariantes $g^{rr} = -\frac{\Delta}{\Sigma}$ y $g^{\theta\theta} = -\frac{1}{\Sigma}$:
\begin{equation}
    g^{tt} E^2 - 2g^{t\phi} E L_z + g^{\phi\phi} L_z^2 - \frac{\Delta}{\Sigma} \left(\frac{dS_r}{dr}\right)^2 - \frac{1}{\Sigma} \left(\frac{dS_\theta}{d\theta}\right)^2 = 0
    \label{eq:HJ_sustituido}
\end{equation}

Puesto que la función de Boyer-Lindquist $\Sigma(r,\theta) = r^2 + a^2 \cos^2\theta > 0$ es estrictamente positiva fuera de la singularidad, \textbf{multiplicamos toda la Ecuación~(\ref{eq:HJ_sustituido}) por $\Sigma(r,\theta)$} para simplificar los denominadores espaciales de las derivadas de la acción:
\begin{equation}
    \Sigma \left( g^{tt} E^2 - 2g^{t\phi} E L_z + g^{\phi\phi} L_z^2 \right) - \Delta \left(\frac{dS_r}{dr}\right)^2 - \left(\frac{dS_\theta}{d\theta}\right)^2 = 0
    \label{eq:HJ_multiplicado_sigma}
\end{equation}

Al reemplazar en el primer término de la Ecuación~(\ref{eq:HJ_multiplicado_sigma}) las componentes contravariantes $g^{tt}, g^{t\phi}, g^{\phi\phi}$ deducidas previamente en las Ecuaciones~(\ref{eq:g_contra_tt}--\ref{eq:g_contra_phiphi}), el factor $\Sigma$ del denominador general se simplifica idénticamente. Agrupando los trinomios cuadrados perfectos en el numerador, se obtiene la descomposición exacta \citep{pdfcoffee.com_chandrasekhar, Grenzebach2014}:
\begin{equation}
    \Sigma \left( g^{tt} E^2 - 2g^{t\phi} E L_z + g^{\phi\phi} L_z^2 \right) = \frac{1}{\Delta(r)} \left[ \frac{E}{c}(r^2+a^2) - a L_z \right]^2 - \left( \frac{a E}{c} \sin\theta - \frac{L_z}{\sin\theta} \right)^2
    \label{eq:desarrollo_binomios_simplificado}
\end{equation}

Sustituyendo la Ecuación~(\ref{eq:desarrollo_binomios_simplificado}) en la Ecuación~(\ref{eq:HJ_multiplicado_sigma}):
\begin{equation}
    \frac{1}{\Delta(r)} \left[ \frac{E}{c}(r^2+a^2) - a L_z \right]^2 - \left( \frac{a E}{c} \sin\theta - \frac{L_z}{\sin\theta} \right)^2 - \Delta(r) \left(\frac{dS_r}{dr}\right)^2 - \left(\frac{dS_\theta}{d\theta}\right)^2 = 0
    \label{eq:HJ_desplegado_completo}
\end{equation}

Reorganizamos la Ecuación~(\ref{eq:HJ_desplegado_completo}) aislando la totalidad de los términos que dependen de la coordenada radial $r$ a la izquierda y los que dependen del ángulo polar $\theta$ a la derecha:
\begin{equation}
    \frac{1}{\Delta(r)} \left[ \frac{E}{c}(r^2+a^2) - a L_z \right]^2 - \Delta(r) \left(\frac{dS_r}{dr}\right)^2 = \left(\frac{dS_\theta}{d\theta}\right)^2 + \left( \frac{a E}{c} \sin\theta - \frac{L_z}{\sin\theta} \right)^2
    \label{eq:separacion_estricta}
\end{equation}

Dado que el lado izquierdo de la Ecuación~(\ref{eq:separacion_estricta}) es una función exclusiva de $r$ y el lado derecho es una función exclusiva de $\theta$, la única forma de que la igualdad se mantenga para todo el espacio-tiempo es que ambas expresiones sean iguales a una constante de separación independiente. Siguiendo la convención clásica de Carter \citep{pdfcoffee.com_chandrasekhar, 1403.5234v2}, definimos dicha constante como $\mathcal{K} + \left(L_z - \frac{a E}{c}\right)^2$, donde $\mathcal{K}$ es la constante de Carter.

\subsection{Formulación explícita del Potencial Radial Efectivo y la velocidad radial}

De la parte radial de la Ecuación~(\ref{eq:separacion_estricta}) igualada a la constante de separación:
\begin{equation}
    \frac{1}{\Delta(r)} \left[ \frac{E}{c}(r^2+a^2) - a L_z \right]^2 - \Delta(r) \left(\frac{dS_r}{dr}\right)^2 = \mathcal{K} + \left(L_z - \frac{a E}{c}\right)^2
    \label{eq:radial_separada_paso1}
\end{equation}

Multiplicando toda la Ecuación~(\ref{eq:radial_separada_paso1}) por la función del horizonte $\Delta(r)$:
\begin{equation}
    \left[ \frac{E}{c}(r^2+a^2) - a L_z \right]^2 - \Delta^2(r) \left(\frac{dS_r}{dr}\right)^2 = \Delta(r) \left[ \mathcal{K} + \left(L_z - \frac{a E}{c}\right)^2 \right]
    \label{eq:radial_separada_paso2}
\end{equation}

Despejando el término derivativo de la acción radial $\Delta^2(r) \left(\frac{dS_r}{dr}\right)^2$, definimos formalmente el \textbf{Potencial Radial Efectivo} $\mathcal{R}(r)$ en unidades del SI ($\text{J}^2\cdot\text{m}^2$) \citep{1403.5234v2, Kobialko2024}:
\begin{equation}
    \mathcal{R}(r) \equiv \Delta^2(r) \left(\frac{dS_r}{dr}\right)^2 = \left[ \frac{E}{c}(r^2+a^2) - a L_z \right]^2 - \Delta(r) \left[ \mathcal{K} + \left(L_z - \frac{a E}{c}\right)^2 \right]
    \label{eq:potencial_R_SI_final}
\end{equation}

Finalmente, conectamos la derivada de la acción radial $\frac{dS_r}{dr}$ con la velocidad coordenada radial $\dot{r} \equiv \frac{dr}{d\lambda}$. Empleando la relación del impulso radial $p_r = \frac{dS_r}{dr} = -\frac{\Sigma}{\Delta} \dot{r}$, tenemos:
\begin{equation}
    \Delta(r) \left(\frac{dS_r}{dr}\right) = -\Sigma(r,\theta) \, \dot{r}
    \label{eq:relacion_p_r_dot_r}
\end{equation}
Elevando la Ecuación~(\ref{eq:relacion_p_r_dot_r}) al cuadrado e igualando con la definición del potencial $\mathcal{R}(r)$ de la Ecuación~(\ref{eq:potencial_R_SI_final}):
\begin{equation}
    \Sigma^2(r,\theta) \left( \frac{dr}{d\lambda} \right)^2 = \mathcal{R}(r) \quad \implies \quad \Sigma(r,\theta) \frac{dr}{d\lambda} = \pm \sqrt{\mathcal{R}(r)}
    \label{eq:ecuacion_movimiento_r_final}
\end{equation}





% ==============================================================================
% SECCIÓN: DEDUCCIÓN DE LA FUNCIÓN R^2(r) DE KOBIALKO Y PARÁMETROS DE IMPACTO (Phi, Q)
% ==============================================================================

\section{Deducción a partir de la función \texorpdfstring{$R^2(r)$}{R^2(r)} de Kobialko y los parámetros de impacto \texorpdfstring{$(\Phi, \mathcal{Q})$}{(\Phi, \mathcal{Q})}}

Para describir la dinámica geodésica de fotones ($m=0$) en el espaciotiempo de Kerr-Newman, el espacio de fases se parametriza mediante las constantes de movimiento conservadas: energía $E$, momento angular azimutal $L_z$ y la constante de Carter $\mathcal{K}$ \citep{1403.5234v2, pdfcoffee.com_chandrasekhar}. Definimos los dos \textbf{parámetros de impacto reducidos} con dimensiones de longitud y área en el Sistema Internacional (SI) como \citep{Grenzebach2014, Kobialko2024}:
\begin{equation}
    \Phi \equiv \frac{c L_z}{E} \quad [\text{m}]
    \label{eq:def_Phi_SI}
\end{equation}
\begin{equation}
    \mathcal{Q} \equiv \frac{c^2 \mathcal{K}}{E^2} \quad [\text{m}^2]
    \label{eq:def_Q_SI}
\end{equation}

Dividiendo el potencial radial efectivo de Hamilton-Jacobi $\mathcal{R}(r)$ entre $E^2/c^2$, obtenemos la función adimensional escalada de Kobialko $R^2(r)$ para la métrica de Kerr-Newman \citep{perturbaciones_en_bh, Kobialko2024}:
\begin{equation}
    R^2(r) \equiv \frac{c^2 \mathcal{R}(r)}{E^2} = \left( r^2 + a^2 - a \Phi \right)^2 - \Delta(r) \left[ \mathcal{Q} + (\Phi - a)^2 \right]
    \label{eq:kobialko_R2_SI_general}
\end{equation}
Sustituyendo la función del horizonte $\Delta(r) = r^2 - 2 r_g r + a^2 + r_g^2 \delta$, donde $r_g \equiv \frac{GM}{c^2}$ es el radio gravitacional y $\delta \equiv \frac{Q^2}{4\pi\varepsilon_0 G M^2} \ll 1$ es el parámetro adimensional de carga eléctrica débil, la función $R^2(r)$ se expande como:
\begin{equation}
    R^2(r) = \left( r^2 + a^2 - a \Phi \right)^2 - \left( r^2 - 2 r_g r + a^2 + r_g^2 \delta \right) \left[ \mathcal{Q} + (\Phi - a)^2 \right]
    \label{eq:kobialko_R2_expandida_general}
\end{equation}

Las condiciones matemáticas analíticas para la existencia de una órbita fotónica esférica inestable (esfera de fotones) en el radio $r_{\mathrm{ph}}$ exigen la anulación simultánea del potencial y de su primera derivada radial \citep{prueba_bh_cargado, dinamica_de_photones_bh_cargado, perturbaciones_en_bh}:
\begin{equation}
    \text{Condición 1:} \quad R^2(r_{\mathrm{ph}}) = 0
    \label{eq:cond1_R2}
\end{equation}
\begin{equation}
    \text{Condición 2:} \quad \left. \frac{d(R^2)}{dr} \right|_{r_{\mathrm{ph}}} = 0
    \label{eq:cond2_R2}
\end{equation}

Al resolver el sistema de ecuaciones~(\ref{eq:cond1_R2})--(\ref{eq:cond2_R2}), se obtienen los valores críticos de los parámetros de impacto $(\Phi_c, \mathcal{Q}_c)$ que parametrizan la frontera de la sombra celeste \citep{1403.5234v2, Grenzebach2014}:
\begin{equation}
    \Phi_c(r_{\mathrm{ph}}) = \frac{(r_{\mathrm{ph}}^2 + a^2) \Delta'(r_{\mathrm{ph}}) - 4 r_{\mathrm{ph}} \Delta(r_{\mathrm{ph}})}{a \, \Delta'(r_{\mathrm{ph}})}
    \label{eq:Phi_critico}
\end{equation}
\begin{equation}
    \mathcal{Q}_c(r_{\mathrm{ph}}) = \frac{16 r_{\mathrm{ph}}^2 \Delta(r_{\mathrm{ph}}) - \left[ \Delta'(r_{\mathrm{ph}}) (r_{\mathrm{ph}}^2 - a^2) - 4 r_{\mathrm{ph}} \Delta(r_{\mathrm{ph}}) \right]^2}{a^2 [\Delta'(r_{\mathrm{ph}})]^2}
    \label{eq:Q_critico}
\end{equation}
donde $\Delta'(r) \equiv \frac{d\Delta}{dr} = 2r - 2r_g$. Se define el \textbf{parámetro de impacto crítico combinado al cuadrado} $b_c^2(r_{\mathrm{ph}})$ mediante la suma cuadrática \citep{perturbaciones_en_bh, Kobialko2024}:
\begin{equation}
    b_c^2(r_{\mathrm{ph}}) \equiv \mathcal{Q}_c(r_{\mathrm{ph}}) + \Phi_c^2(r_{\mathrm{ph}}) = \frac{16 r_{\mathrm{ph}}^2 \Delta(r_{\mathrm{ph}})}{[\Delta'(r_{\mathrm{ph}})]^2}
    \label{eq:bc2_combinado_def}
\end{equation}

\subsection{Deducción detallada en el límite estático y expansión de Kobialko}

En el límite esféricamente simétrico sin rotación ($a \to 0$, métrica de Reissner-Nordström), la función de horizonte es $\Delta(r) = r^2 - 2 r_g r + r_g^2 \delta$ y su derivada resulta $\Delta'(r) = 2(r - r_g)$. Sustituyendo esta derivada en la ec.~(\ref{eq:bc2_combinado_def}):
\begin{equation}
    b_c^2(r_{\mathrm{ph}}) = \frac{16 r_{\mathrm{ph}}^2 \Delta(r_{\mathrm{ph}})}{[2(r_{\mathrm{ph}} - r_g)]^2} = \frac{4 r_{\mathrm{ph}}^2 \Delta(r_{\mathrm{ph}})}{(r_{\mathrm{ph}} - r_g)^2}
    \label{eq:bc2_limite_estatico_1}
\end{equation}
Por otro lado, de la Condición 1 ec.~(\ref{eq:cond1_R2}) con $a=0$, tenemos $R^2(r_{\mathrm{ph}}) = r_{\mathrm{ph}}^4 - \Delta(r_{\mathrm{ph}}) b_c^2 = 0$, de donde:
\begin{equation}
    b_c^2(r_{\mathrm{ph}}) = \frac{r_{\mathrm{ph}}^4}{\Delta(r_{\mathrm{ph}})}
    \label{eq:bc2_limite_estatico_2}
\end{equation}

Igualando la ec.~(\ref{eq:bc2_limite_estatico_1}) con la ec.~(\ref{eq:bc2_limite_estatico_2}):
\begin{equation}
    \frac{r_{\mathrm{ph}}^4}{\Delta(r_{\mathrm{ph}})} = \frac{4 r_{\mathrm{ph}}^2 \Delta(r_{\mathrm{ph}})}{(r_{\mathrm{ph}} - r_g)^2}
    \label{eq:igualdad_bc2}
\end{equation}
Simplificando $r_{\mathrm{ph}}^2 \neq 0$ y tomando la raíz cuadrada:
\begin{equation}
    r_{\mathrm{ph}} (r_{\mathrm{ph}} - r_g) = 2 \Delta(r_{\mathrm{ph}})
    \label{eq:ec_rph_paso1}
\end{equation}
Sustituyendo la expresión explícita de $\Delta(r_{\mathrm{ph}}) = r_{\mathrm{ph}}^2 - 2 r_g r_{\text{ph}} + r_g^2 \delta$:
\begin{equation}
    r_{\mathrm{ph}}^2 - r_g r_{\mathrm{ph}} = 2 \left( r_{\mathrm{ph}}^2 - 2 r_g r_{\mathrm{ph}} + r_g^2 \delta \right)
    \label{eq:ec_rph_paso2}
\end{equation}
Reagrupando los términos algebraicos se obtiene la ecuación cuadrática fundamental para el radio de la esfera de fotones:
\begin{equation}
    r_{\mathrm{ph}}^2 - 3 r_g r_{\mathrm{ph}} + 2 r_g^2 \delta = 0
    \label{eq:cuadratica_rph}
\end{equation}

Resolviendo de forma exacta la ec.~(\ref{eq:cuadratica_rph}) mediante la fórmula binomial para la raíz inestable exterior:
\begin{equation}
    r_{\mathrm{ph}}(\delta) = \frac{3 r_g + \sqrt{9 r_g^2 - 8 r_g^2 \delta}}{2} = \frac{3 r_g}{2} \left( 1 + \sqrt{1 - \frac{8}{9}\delta} \right)
    \label{eq:rph_exacto}
\end{equation}
Realizamos la expansión en serie de Taylor de la función binomia $\sqrt{1 - x} = 1 - \frac{1}{2}x - \frac{1}{8}x^2 - \dots$ tomando $x = \frac{8}{9}\delta \ll 1$ \citep{prueba_bh_cargado}:
\begin{equation}
    \sqrt{1 - \frac{8}{9}\delta} = 1 - \frac{1}{2}\left(\frac{8}{9}\delta\right) - \frac{1}{8}\left(\frac{8}{9}\delta\right)^2 - \mathcal{O}(\delta^3)
    \label{eq:taylor_paso1}
\end{equation}
\begin{equation}
    \sqrt{1 - \frac{8}{9}\delta} = 1 - \frac{4}{9}\delta - \frac{8}{81}\delta^2 - \mathcal{O}(\delta^3)
    \label{eq:taylor_raiz}
\end{equation}
Sustituyendo la ec.~(\ref{eq:taylor_raiz}) en la ec.~(\ref{eq:rph_exacto}):
\begin{equation}
    r_{\mathrm{ph}}(\delta) = \frac{3 r_g}{2} \left( 2 - \frac{4}{9}\delta - \frac{8}{81}\delta^2 \right) = 3 r_g \left( 1 - \frac{2}{9}\delta - \frac{4}{81}\delta^2 + \mathcal{O}(\delta^3) \right)
    \label{eq:rph_perturbado_SI}
\end{equation}

Para calcular la expansión del parámetro de impacto crítico al cuadrado $b_c^2(\delta)$, notamos que de la ec.~(\ref{eq:cuadratica_rph}) se cumple la identidad $\Delta(r_{\mathrm{ph}}) = r_g (r_{\mathrm{ph}} - r_g \delta)$. Sustituyendo esto en la ec.~(\ref{eq:bc2_limite_estatico_2}):
\begin{equation}
    b_c^2(\delta) = \frac{r_{\mathrm{ph}}^3(\delta)}{r_g \left( 1 - \frac{r_g \delta}{r_{\mathrm{ph}}} \right)}
    \label{eq:bc2_relacion_simplificada}
\end{equation}

Elevando la ec.~(\ref{eq:rph_perturbado_SI}) al cubo:
\begin{equation}
    r_{\mathrm{ph}}^3(\delta) = 27 r_g^3 \left( 1 - \frac{2}{9}\delta - \frac{4}{81}\delta^2 \right)^3 = 27 r_g^3 \left( 1 - \frac{2}{3}\delta + \mathcal{O}(\delta^3) \right)
    \label{eq:rph_cubo}
\end{equation}
Por su parte, el denominador se expande como:
\begin{equation}
    1 - \frac{r_g \delta}{r_{\mathrm{ph}}} = 1 - \frac{\delta}{3\left(1 - \frac{2}{9}\delta\right)} = 1 - \frac{1}{3}\delta - \frac{2}{27}\delta^2 + \mathcal{O}(\delta^3)
    \label{eq:denominador_exp}
\end{equation}

Dividiendo la ec.~(\ref{eq:rph_cubo}) entre la ec.~(\ref{eq:denominador_exp}):
\begin{align}
    b_c^2(\delta) &= \frac{27 r_g^2 \left( 1 - \frac{2}{3}\delta \right)}{1 - \frac{1}{3}\delta - \frac{2}{27}\delta^2} = 27 r_g^2 \left( 1 - \frac{2}{3}\delta \right) \left( 1 + \frac{1}{3}\delta + \frac{5}{27}\delta^2 \right) \nonumber \\[0.2cm]
    &= 27 r_g^2 \left[ 1 + \left( \frac{1}{3} - \frac{2}{3} \right)\delta + \left( \frac{5}{27} - \frac{2}{9} \right)\delta^2 + \mathcal{O}(\delta^3) \right]
\end{align}
Obteniendo finalmente la célebre **expansión analítica de Kobialko** para el parámetro de impacto crítico al cuadrado en unidades SI \citep{perturbaciones_en_bh, Kobialko2024}:
\begin{equation}
    b_c^2(\delta) = 27 r_g^2 \left( 1 - \frac{1}{3}\delta - \frac{1}{27}\delta^2 + \mathcal{O}(\delta^3) \right)
    \label{eq:bc2_kobialko_SI}
\end{equation}
donde el coeficiente de segundo orden $D = 1/27$ constituye la **firma unívoca de electrovacío** del espaciotiempo cargado frente a soluciones alternativas de gravedad modificada.
 %% 3o capítulo
%=============================
%  CUARTO CAPITULO
%=============================
\chapter{ANÁLISIS PERTURBATIVO EN KERR-NEWMAN Y ESTABILIDAD DE LYAPUNOV}

\section{Potencial radial efectivo \texorpdfstring{$\mathcal{R}(r)$}{R(r)} en el plano ecuatorial}
Consideramos el espaciotiempo estacionario y axialmente simétrico de Kerr-Newman en coordenadas de Boyer-Lindquist $(t, r, \theta, \phi)$ expresado en el Sistema Internacional de Unidades (SI) con la signatura métrica de Minkowski $(+,-,-,-)$ \citep{Carroll, Hobson}:
\begin{equation}
    \Delta(r) = r^2 - 2 r_g r + a^2 + r_Q^2
    \label{eq:delta_KN_SI}
\end{equation}
donde $r_g \equiv \frac{GM}{c^2}$ es el radio gravitacional, $a \equiv \frac{J}{Mc}$ es el parámetro de espín por unidad de masa y $r_Q^2 \equiv \frac{G Q^2}{4\pi\varepsilon_0 c^4} = 2 r_g^2 \delta$, siendo $\delta \equiv \frac{Q^2}{4\pi\varepsilon_0 G M^2} \ll 1$ el parámetro adimensional de carga eléctrica débil \citep{prueba_bh_cargado}.

Para geodésicas nulas ($ds^2 = 0$) restringidas al plano ecuatorial ($\theta = \pi/2$, $p_\theta = 0$), la constante de separación de Carter adopta la forma $\mathcal{K} = (L_z - a E/c)^2$ \citep{Chandrasekhar, 1403.5234v2}. La función de la acción de Hamilton-Jacobi $S = -E t + L_z \phi + S_r(r)$ conduce a la ecuación diferencial radial exacta \citep{Carroll, Hobson}:
\begin{equation}
    r^4 \left( \frac{dr}{d\lambda} \right)^2 = \mathcal{R}(r)
    \label{eq:partida_KN_R}
\end{equation}
donde $\lambda$ es el parámetro afín del fotón y el potencial radial efectivo $\mathcal{R}(r)$ está dado explícitamente por \citep{Chandrasekhar, 1403.5234v2}:
\begin{equation}
    \mathcal{R}(r) = \left[ E (r^2 + a^2) - a L_z c \right]^2 - \Delta(r) \left[ (L_z c - a E)^2 \right]
    \label{eq:potencial_carter_KN}
\end{equation}

Definiendo el parámetro de impacto ecuatorial $b \equiv \frac{L_z c}{E}$ (con dimensión de longitud), reescribimos el potencial radial normalizado por $E^2$ \citep{Sneppen, chaos_fractals_photon_ring}:
\begin{equation}
    \frac{\mathcal{R}(r)}{E^2} = \left( r^2 + a^2 - a b \right)^2 - \left( r^2 - 2 r_g r + a^2 + 2 r_g^2 \delta \right) (b - a)^2
    \label{eq:potencial_R_parametro_b}
\end{equation}

\section{Condiciones de equilibrio y expansión perturbativa en carga débil}
Las órbitas circulares inestables de fotones que forman la esfera (o banda de lentes) de fotones de radio crítico $r_{ph}$ en el plano ecuatorial cumplen las dos condiciones extremales de equilibrio \citep{fhoton_bh, prueba_bh_cargado}:
\begin{equation}
    \mathcal{R}(r_{ph}) = 0, \quad \left. \frac{d\mathcal{R}}{dr} \right|_{r_{ph}} = 0
    \label{eq:condiciones_equilibrio_KN}
\end{equation}

En el régimen de carga eléctrica débil ($\delta \ll 1$), proponemos una expansión en serie de potencias regulares para el radio crítico $r_{ph}$ y para el parámetro de impacto crítico $b_c$ alrededor de la solución no perturbada de Kerr ($\delta = 0$) \citep{Ali-H.-Nayfeh-Introduction-to-perturbation-techniques-1993-Wiley-VCH, TFG-G7513}:
\begin{align}
    r_{ph}(a, \delta) &= r_{ph}^{(0)}(a) + \delta \, r_{ph}^{(1)}(a) + \delta^2 \, r_{ph}^{(2)}(a) + \mathcal{O}(\delta^3) \label{eq:expansion_rph} \\
    b_c(a, \delta) &= b_c^{(0)}(a) + \delta \, b_c^{(1)}(a) + \delta^2 \, b_c^{(2)}(a) + \mathcal{O}(\delta^3) \label{eq:expansion_bc}
\end{align}

El radio de fotones no perturbado de Kerr $r_{ph}^{(0)}(a)$ para órbitas ecuatoriales progradas ($b > 0$) y retrógradas ($b < 0$) viene dado por la expresión analítica cerrada \citep{Hobson, Achour_2025_J._Cosmol._Astropart._Phys._2025_012}:
\begin{equation}
    r_{ph}^{(0)}(a) = 2 r_g \left\{ 1 + \cos \left[ \frac{2}{3} \arccos \left( \mp \frac{a}{r_g} \right) \right] \right\}
    \label{eq:rph_kerr_exacto}
\end{equation}

Sustituyendo las expansiones~(\ref{eq:expansion_rph}) y~(\ref{eq:expansion_bc}) en las condiciones de equilibrio~(\ref{eq:condiciones_equilibrio_KN}) e igualando términos de orden $\mathcal{O}(\delta^1)$, obtenemos la corrección de primer orden al radio crítico provocada por la carga:
\begin{equation}
    r_{ph}^{(1)}(a) = -\frac{2 r_g^2 \left( b_c^{(0)} - a \right)^2}{4 r_{ph}^{(0)} \left( r_{ph}^{(0)2} + a^2 - a b_c^{(0)} \right) - 2 \left( r_{ph}^{(0)} - r_g \right) \left( b_c^{(0)} - a \right)^2} < 0
    \label{eq:rph_primer_orden_KN}
\end{equation}
La Ecuación~(\ref{eq:rph_primer_orden_KN}) demuestra rigurosamente que la presencia de una carga eléctrica residual $Q$ disminuye el radio de la esfera de fotones respecto al caso neutro de Kerr ($r_{ph} < r_{ph}^{(0)}$), debido a la atracción electrogravitacional adicional \citep{Heydarzade_2024}.

\section{Deducción del exponente de Lyapunov y oscilador inestable}
Para determinar la tasa de divergencia de las trayectorias fotónicas vecinas, introducimos una pequeña perturbación radial $\epsilon(\lambda) \equiv r(\lambda) - r_{ph}$ alrededor de la órbita de equilibrio inestable \citep{Sneppen, chaos_fractals_photon_ring}. Expandimos el potencial radial $\mathcal{R}(r)$ en serie de Taylor hasta segundo orden:
\begin{equation}
    \mathcal{R}(r) \approx \mathcal{R}(r_{ph}) + \left.\frac{d\mathcal{R}}{dr}\right|_{r_{ph}} \epsilon + \frac{1}{2} \left.\frac{d^2\mathcal{R}}{dr^2}\right|_{r_{ph}} \epsilon^2 = \frac{1}{2} \left.\frac{d^2\mathcal{R}}{dr^2}\right|_{r_{ph}} \epsilon^2
    \label{eq:R_segundo_orden_KN}
\end{equation}

Sustituyendo la Ecuación~(\ref{eq:R_segundo_orden_KN}) en la ecuación de movimiento radial~(\ref{eq:partida_KN_R}) y tomando la raíz cuadrada:
\begin{equation}
    \frac{d\epsilon}{d\lambda} = \pm \frac{\sqrt{\left.\frac{d^2\mathcal{R}}{dr^2}\right|_{r_{ph}}}}{\sqrt{2} r_{ph}^2} \epsilon(\lambda)
    \label{eq:d_epsilon_lambda_KN}
\end{equation}

Derivando respecto al parámetro afín $\lambda$, obtenemos la ecuación diferencial del **oscilador armónico inestable de Lyapunov** \citep{Sneppen, 2406.09498v1}:
\begin{equation}
    \frac{d^2\epsilon}{d\lambda^2} - \gamma_\lambda^2 \epsilon(\lambda) = 0
    \label{eq:oscilador_lyapunov_SI_KN}
\end{equation}
donde el exponente de Lyapunov afín $\gamma_\lambda$ está definido por:
\begin{equation}
    \gamma_\lambda(a, \delta) \equiv \frac{\sqrt{\left.\frac{d^2\mathcal{R}}{dr^2}\right|_{r_{ph}}}}{\sqrt{2} r_{ph}^2}
    \label{eq:def_gamma_lambda_KN}
\end{equation}

\subsection{Transformación al exponente orbital \texorpdfstring{$\gamma_\phi$}{gamma\_phi}}
Debido al acoplamiento de rotación en la métrica no diagonal de Kerr-Newman, la velocidad angular azimutal en el radio de la esfera de fotones $\left.\frac{d\phi}{d\lambda}\right|_{r_{ph}}$ se obtiene de las ecuaciones geodésicas \citep{Carroll, Chandrasekhar}:
\begin{equation}
    \left.\frac{d\phi}{d\lambda}\right|_{r_{ph}} = \frac{E}{r_{ph}^2 \Delta(r_{ph})} \left[ a \left( r_{ph}^2 + a^2 - a b_c \right) + \Delta(r_{ph}) \left( b_c - a \right) \right]
    \label{eq:dphi_dlambda_KN}
\end{equation}

Aplicando la regla de la cadena $\frac{d\epsilon}{d\lambda} = \frac{d\epsilon}{d\phi} \left( \frac{d\phi}{d\lambda} \right)$, la EDO orbital toma la forma \citep{Sneppen}:
\begin{equation}
    \frac{d^2\epsilon}{d\phi^2} - \gamma_\phi^2 \epsilon(\phi) = 0
    \label{eq:oscilador_orbital_phi_KN}
\end{equation}
donde el **exponente de Lyapunov orbital por radián** $\gamma_\phi(a, \delta)$ en el espaciotiempo de Kerr-Newman se expresa analíticamente como:
\begin{equation}
    \gamma_\phi(a, \delta) \equiv \frac{\gamma_\lambda}{\left.\frac{d\phi}{d\lambda}\right|_{r_{ph}}} = \frac{\Delta(r_{ph}) \sqrt{\left.\frac{d^2\mathcal{R}}{dr^2}\right|_{r_{ph}}}}{\sqrt{2} E \left[ a \left( r_{ph}^2 + a^2 - a b_c \right) + \Delta(r_{ph}) \left( b_c - a \right) \right]}
    \label{eq:def_gamma_phi_KN_exacto}
\end{equation}

Para fotones progrados con alto espín ($a \to r_g$), el arrastre de marcos reduce el exponente de Lyapunov orbital a $\gamma_\phi < 1$, lo que atenúa la divergencia radial y atenúa la atenuación de Sneppen ($f_{thin} = e^{-2\pi\gamma_\phi} \gg e^{-2\pi}$), fundamentando el ensanchamiento del subanillo $n=1$ y su detectabilidad interferométrica por el ngEHT y BHEX \citep{2406.09498v1}.

\section{Deducción paso a paso del Exponente de Lyapunov \texorpdfstring{$\gamma$}{gamma}}
Deducimos ahora de manera rigurosa la ecuación del exponente de Lyapunov a partir del potencial radial.
Partimos de la Ecuación~(\ref{eq:eq_radial_R}):
\begin{equation}
    r^4 \left(\frac{dr}{d\lambda}\right)^2 = \mathcal{R}(r)
    \label{eq:partida_lyapunov}
\end{equation}
Realizamos el desarrollo en serie de Taylor del potencial $\mathcal{R}(r)$ alrededor del radio crítico $r_{ph}$:
\begin{equation}
    \mathcal{R}(r) = \mathcal{R}(r_{ph}) + \left.\frac{d\mathcal{R}}{dr}\right|_{r_{ph}} (r - r_{ph}) + \frac{1}{2}\left.\frac{d^2\mathcal{R}}{dr^2}\right|_{r_{ph}} (r - r_{ph})^2 + \mathcal{O}\left((r - r_{ph})^3\right)
    \label{eq:taylor_R}
\end{equation}
Sustituyendo la Ecuación~(\ref{eq:cond1_R2}) ($\mathcal{R}(r_{ph}) = 0$) y la Ecuación~(\ref{eq:cond2_R2}) ($\left.\frac{d\mathcal{R}}{dr}\right|_{r_{ph}} = 0$) en la Ecuación~(\ref{eq:taylor_R}), los dos primeros términos se anulan exactamente, dejando:
\begin{equation}
    \mathcal{R}(r) \approx \frac{1}{2}\left.\frac{d^2\mathcal{R}}{dr^2}\right|_{r_{ph}} (r - r_{ph})^2
    \label{eq:R_segundo_orden}
\end{equation}
Sustituyendo la Ecuación~(\ref{eq:R_segundo_orden}) en la Ecuación~(\ref{eq:partida_lyapunov}):
\begin{equation}
    r_{ph}^4 \left(\frac{dr}{d\lambda}\right)^2 = \frac{1}{2}\left.\frac{d^2\mathcal{R}}{dr^2}\right|_{r_{ph}} (r - r_{ph})^2
    \label{eq:radial_linealizada}
\end{equation}
Definimos la pequeña desviación radial $\epsilon(\lambda) \equiv r(\lambda) - r_{ph}$. Tomando la raíz cuadrada en la Ecuación~(\ref{eq:radial_linealizada}):
\begin{equation}
    \frac{d\epsilon}{d\lambda} = \pm \frac{\sqrt{\left.\frac{d^2\mathcal{R}}{dr^2}\right|_{r_{ph}}}}{\sqrt{2} r_{ph}^2} \epsilon(\lambda)
    \label{eq:d_epsilon_lambda}
\end{equation}
Derivando la Ecuación~(\ref{eq:d_epsilon_lambda}) respecto al parámetro afín $\lambda$, obtenemos la ecuación del oscilador inestable paso a paso:
\begin{equation}
    \frac{d^2\epsilon}{d\lambda^2} - \gamma_\lambda^2 \epsilon(\lambda) = 0
    \label{eq:oscilador_inestable_SI}
\end{equation}
donde el exponente de Lyapunov respecto al parámetro afín se define exactamente como:
\begin{equation}
    \gamma_\lambda \equiv \frac{\sqrt{\left.\frac{d^2\mathcal{R}}{dr^2}\right|_{r_{ph}}}}{\sqrt{2} r_{ph}^2}
    \label{eq:def_lyapunov_SI}
\end{equation}
La solución general de la Ecuación~(\ref{eq:oscilador_inestable_SI}) es:
\begin{equation}
    \epsilon(\lambda) = C_1 e^{\gamma_\lambda \lambda} + C_2 e^{-\gamma_\lambda \lambda}
    \label{eq:solucion_lyapunov}
\end{equation}
Esto demuestra analíticamente que la presencia de la raíz exponencial $e^{\gamma_\lambda \lambda}$ en la Ecuación~(\ref{eq:solucion_lyapunov}) impulsa la divergencia de las trayectorias vecinas.
 %% 4o capítulo
%=============================
%  QUINTO CAPITULO
%=============================
\chapter{ESTABILIDAD, CAOS Y LA JERARQUÍA DE ANILLOS DE FOTONES}

\section{El factor de adelgazamiento de Sneppen \texorpdfstring{$f_{thin} = e^{-2\pi\gamma}$}{f\_thin = e^{-2 pi gamma}}}
\citet{Sneppen} demostró que los fotones atrapados momentáneamente cerca de la esfera de fotones forman una jerarquía fractal de sub-anillos \citep{Sneppen, pertbuciiones_teoria_en_bh}. Cada sub-anillo $n$ corresponde a fotones que completan $n$ medias órbitas. La relación de espesores entre sub-anillos consecutivos se rige por el factor de adelgazamiento de Sneppen \citep{Sneppen}:
\begin{equation}
    f_{thin} = e^{-2\pi\gamma_\phi}
    \label{eq:factor_sneppen}
\end{equation}
En el caso estático de Schwarzschild, $\gamma_\phi = 1$, lo que da $f_{thin} = e^{-2\pi} \approx 1/535.5$ \citep{Sneppen, chaos_fractals_photon_ring}.

\section{Ensanchamiento de los anillos por rotación y carga}
En el espaciotiempo de Kerr-Newman, la rotación degrada el exponente de Lyapunov para órbitas progradas ($\gamma_\phi < 1$), reduciendo el factor de adelgazamiento de 535 a valores tan bajos como $f_{thin} \approx 2$ \citep{Sneppen, bh_mergers_photon_ring}. Esto ensancha los sub-anillos $n \ge 1$, permitiendo que sean detectados por instrumentos como el ngEHT y BHEX \citep{Measuring_Photon_Rings_with_the_ngEHT, 2406.09498v1}.

\section{Transición al caos en el espacio de fases}
Mientras que Kerr-Newman es integrable debido a la constante de Carter \citep{1403.5234v2, Achour_2025_J._Cosmol._Astropart._Phys._2025_012}, la presencia de perturbaciones no integrables destruye los toros de Kolmogorov-Arnold-Moser (KAM), generando caos determinista en el espacio de fases \citep{chaos_fractals_photon_ring}.
\label{cap5}
 %% 5o capítulo
\chapter{FIRMAS DE SEGUNDO ORDEN DE KOBIALKO Y DISCRIMINACIÓN DE MODELOS}

\section{Expansión de Taylor de Kobialko para la sombra \texorpdfstring{$R^2(\delta)$}{R^2(delta)}}
Siguiendo a Kobialko \& Gal'tsov (2024) \citep{perturbaciones_en_bh}, expandimos el cuadrado del parámetro de impacto crítico $b_c^2$ hasta segundo orden en el parámetro de carga débil $\delta$:
\begin{equation}
    b_c^2(\delta) = 27 r_g^2 \left( 1 - B \delta - D \delta^2 + \mathcal{O}(\delta^3) \right)
    \label{eq:kobialko_expansion_SI}
\end{equation}
donde $r_g = \frac{GM}{c^2}$. Para la métrica de Kerr-Newman, el coeficiente lineal es $B = 1/3$, idéntico al de otros modelos cargados \citep{perturbaciones_en_bh}.

\section{Aislamiento del coeficiente de segundo orden \texorpdfstring{$D = 1/27$}{D = 1/27}}
Deducimos paso a paso el coeficiente de segundo orden para Kerr-Newman cargado débilmente \citep{perturbaciones_en_bh}:
\begin{equation}
    D_{KN} = \frac{1}{27}
\end{equation}
Comparando analíticamente con modelos alternativos \citep{perturbaciones_en_bh}:
\begin{equation}
    D_{Bardeen} = \frac{2}{27}, \quad D_{Hayward} = \frac{8}{243}
\end{equation}
El coeficiente $D = 1/27$ es la firma única e inconfundible de la carga eléctrica de electrovacío en segundo orden.

\section{El observador geodésico y aberración relativista}
El marco local del observador geodésico introduce correcciones por aberración relativista en las coordenadas celestes $(\alpha, \beta)$ \citep{shadows_quintessence_bh, Achour_2025_J._Cosmol._Astropart._Phys._2025_012}, transformando el perfil angular medido.
 %% 6o capítulo
\chapter{RESULTADOS NUMÉRICOS Y CONTRASTE EXPERIMENTAL CON EL EHT}

\section{Simulación geodésica mediante Runge-Kutta de 4to orden (RK4)}
Se implementa un algoritmo en Python para integrar numéricamente las geodésicas nulas usando RK4 \citep{Geod_sicas_en_la_m_trica_de_Schwarzschild__una_soluci_n_num_rica, TFG-G7513}, resolviendo las ecuaciones de movimiento derivadas en el Capítulo III.

\section{Coordenadas celestes \texorpdfstring{$(\alpha, \beta)$}{(alpha, beta)} e intervalo angular en \texorpdfstring{$\mu\text{as}$}{microsegundos de arco}}
Las coordenadas celestes en el plano del observador a distancia $D$ e inclinación $\theta_o$ son \citep{1403.5234v2, 2406.09498v1}:
\begin{equation}
    \alpha = -\frac{b_c}{\sin\theta_o}, \quad \beta = \pm \sqrt{q_c + a^2 \cos^2\theta_o - b_c^2 \cot^2\theta_o}
\end{equation}
El diámetro angular observado en microsegundos de arco ($\mu as$) viene dado por \citep{Broderick_2022_ApJ_927_6, Reproducibility_of_the_First_Image_of_a_Black_Hole_in_the_Galaxy_M87_From_the_Event_Horizon_Telescope_Collaboration}:
\begin{equation}
    \theta_{shadow} = \frac{2 b_c GM}{c^2 D} \text{ rad}
\end{equation}

\section{Contraste con los datos del EHT para M87* y Sgr A*}
Para M87* ($M \approx 6.5 \times 10^9 M_\odot$, $D \approx 16.8 \text{ Mpc}$), $1 r_g \approx 3.81 \, \mu as$ \citep{Broderick_2022_ApJ_927_6, EHT_Paper1}. El EHT midió un diámetro de $42 \pm 3 \, \mu as$ \citep{EHT_Paper1}. La diferencia de $0.74 \, \mu as$ predicha por el término $D\delta^2$ estará al alcance de las líneas de base espaciales del ngEHT \citep{Measuring_Photon_Rings_with_the_ngEHT}.
 %% 7o capítulo
\chapter{CONCLUSIONES}

\begin{enumerate}
    \item Se formuló rigurosamente la métrica y geodésicas de Kerr-Newman con carga débil en unidades SI y firma $(+,-,-,-)$, confirmando que la carga disminuye el radio crítico de la esfera de fotones.
    \item Se derivó paso a paso la ecuación del oscilador inestable para el exponente de Lyapunov $\gamma_\lambda$, justificando la divergencia exponencial de las trayectorias vecinas.
    \item Se demostró que el coeficiente de segundo orden de Kobialko $D = 1/27$ constituye la huella analítica única para distinguir la carga de Kerr-Newman frente a soluciones regulares de Bardeen y Hayward.
    \item Se comprobó que el ensanchamiento de Sneppen en métricas rotantes incrementa la visibilidad interferométrica de los sub-anillos $n \ge 1$ para telescopios de nueva generación.
\end{enumerate}
 %% conclusiones capítulo

\newpage
% =========================================================================
% REFERENCIAS BIBLIOGRÁFICAS (APA 7.ª Edición - FCF UNMSM)
% =========================================================================
\newpage
\chapter*{Referencias Bibliográficas}
\addcontentsline{toc}{chapter}{Referencias Bibliográficas}

\begin{thebibliography}{99}

\bibitem[Alfaro Sotil(2025)]{prueba_bh_cargado} Alfaro Sotil, H. P. (2025). \textit{Análisis perturbativo de las órbitas de fotones y la estabilidad de la esfera de fotones en agujeros negros débilmente cargados} [Informe de Avance de Tesis, Universidad Nacional Mayor de San Marcos]. Repositorio Institucional UNMSM.

\bibitem[Ben Achour et al.(2025)]{Achour_2025_J._Cosmol._Astropart._Phys._2025_012} Ben Achour, J., Gourgoulhon, É., y Roussille, H. (2025). Black hole photon ring beyond General Relativity: An integrable parametrization. \textit{Journal of Cosmology and Astroparticle Physics}, \textit{2025}(10), Artículo 012. https://doi.org/10.1088/1475-7516/2025/10/012

\bibitem[Berens et al.(2024)]{chaos_fractals_photon_ring} Berens, R., Galison, P., Gravely, T., Lupsasca, A., y Stein, L. C. (2024). Chaos and fractals of the black hole photon ring. \textit{Classical and Quantum Gravity}, \textit{41}(10), Artículo 105001. https://doi.org/10.1088/1361-6382/ad333e
\bibitem[Berens et al.(2024)]{Berens2024} Berens, R., Galison, P., Gravely, T., Lupsasca, A., y Stein, L. C. (2024). Chaos and fractals of the black hole photon ring. \textit{Classical and Quantum Gravity}, \textit{41}(10), Artículo 105001. https://doi.org/10.1088/1361-6382/ad333e

\bibitem[Broderick et al.(2022)]{Broderick_2022_ApJ_927_6} Broderick, A. E., Tiede, P., Pesce, D. W., y Gold, R. (2022). Measuring photon rings with the Event Horizon Telescope. \textit{The Astrophysical Journal}, \textit{927}(1), Artículo 6. https://doi.org/10.3847/1538-4357/ac426a

\bibitem[Carnero Morán(2025)]{TFG-G7513} Carnero Morán, P. (2025). \textit{Métodos de perturbación} [Trabajo de Fin de Grado, Universidad de Valladolid]. Repositorio Documental UVa.
\bibitem[Carnero Morán(2025)]{CarneroMoran2025} Carnero Morán, P. (2025). \textit{Métodos de perturbación} [Trabajo de Fin de Grado, Universidad de Valladolid]. Repositorio Documental UVa.

\bibitem[Carroll(2003)]{Sean_Carroll_Spacetime} Carroll, S. M. (2003). \textit{Spacetime and geometry: An introduction to General Relativity}. Pearson.

\bibitem[Chakraborty y Sarkar(2019)]{fhoton_bh} Chakraborty, S., y Sarkar, S. (2019). Photon sphere evolution in dynamical black hole spacetimes. \textit{Physical Review D}, \textit{100}(6), Artículo 064010. https://doi.org/10.1103/PhysRevD.100.064010

\bibitem[Chandrasekhar(1998)]{pdfcoffee.com_chandrasekhar} Chandrasekhar, S. (1998). \textit{The mathematical theory of black holes}. Oxford University Press.
\bibitem[Chandrasekhar(1983)]{Chandrasekhar1983} Chandrasekhar, S. (1983). \textit{The mathematical theory of black holes}. Oxford University Press.

\bibitem[Dias et al.(2015)]{PhysRevLett.114.151101} Dias, O. J. C., Godazgar, M., y Santos, J. E. (2015). Linear stability of Kerr-Newman black holes. \textit{Physical Review Letters}, \textit{114}(15), Artículo 151101. https://doi.org/10.1103/PhysRevLett.114.151101

\bibitem[Duarte-Ortiz et al.(2024)]{Geod_sicas_en_la_m_trica_de_Schwarzschild__una_soluci_n_num_rica} Duarte-Ortiz, C., Leija, S., y Flores, F. (2024). Geodésicas en la métrica de Schwarzschild: Una solución numérica. \textit{Revista Mexicana de Física}, \textit{70}(2), Artículo 021401.

\bibitem[Event Horizon Telescope Collaboration et al.(2019a)]{EHT_Paper1} Event Horizon Telescope Collaboration et al. (2019a). First M87 Event Horizon Telescope results. I. The shadow of the supermassive black hole. \textit{The Astrophysical Journal Letters}, \textit{875}(1), L1. https://doi.org/10.3847/2041-8213/ab0ec7

\bibitem[Event Horizon Telescope Collaboration et al.(2019b)]{EHT_Paper4} Event Horizon Telescope Collaboration et al. (2019b). First M87 Event Horizon Telescope results. IV. Imaging the central supermassive black hole. \textit{The Astrophysical Journal Letters}, \textit{875}(1), L4. https://doi.org/10.3847/2041-8213/ab0e85

\bibitem[Event Horizon Telescope Collaboration et al.(2019c)]{EHT_Paper6} Event Horizon Telescope Collaboration et al. (2019c). First M87 Event Horizon Telescope results. VI. The shadow and mass of the central black hole. \textit{The Astrophysical Journal Letters}, \textit{875}(1), L6. https://doi.org/10.3847/2041-8213/ab1141

\bibitem[Event Horizon Telescope Collaboration et al.(2022)]{EHT_SgrA_Paper1} Event Horizon Telescope Collaboration et al. (2022). First Sagittarius A* Event Horizon Telescope results. I. The shadow of the supermassive black hole. \textit{The Astrophysical Journal Letters}, \textit{930}(1), L12. https://doi.org/10.3847/2041-8213/ac6674

\bibitem[Event Horizon Telescope Collaboration et al.(2024a)]{EHT_SgrA_Paper7} Event Horizon Telescope Collaboration et al. (2024a). First Sagittarius A* Event Horizon Telescope results. VII. Polarization of the supermassive black hole. \textit{The Astrophysical Journal Letters}, \textit{964}(2), L25. https://doi.org/10.3847/2041-8213/ad2df0

\bibitem[Event Horizon Telescope Collaboration et al.(2024b)]{EHT_SgrA_Paper8} Event Horizon Telescope Collaboration et al. (2024b). First Sagittarius A* Event Horizon Telescope results. VIII. Physical implications of the polarimetric properties. \textit{The Astrophysical Journal Letters}, \textit{964}(2), L26. https://doi.org/10.3847/2041-8213/ad2df1

\bibitem[Giataganas et al.(2024)]{bh_mergers_photon_ring} Giataganas, D., Kehagias, A., y Riotto, A. (2024). Black hole mergers as the fastest photon ring scramblers. \textit{Physical Review Letters}, \textit{133}(8), Artículo 081401. https://doi.org/10.1103/PhysRevLett.133.081401

\bibitem[Grenzebach et al.(2014)]{1403.5234v2} Grenzebach, A., Perlick, V., y Lämmerzahl, C. (2014). Photon regions and shadows of Kerr-Newman-NUT black holes with a cosmological constant. \textit{Physical Review D}, \textit{89}(12), Artículo 124004. https://doi.org/10.1103/PhysRevD.89.124004
\bibitem[Grenzebach et al.(2014)]{Grenzebach2014} Grenzebach, A., Perlick, V., y Lämmerzahl, C. (2014). Photon regions and shadows of Kerr-Newman-NUT black holes with a cosmological constant. \textit{Physical Review D}, \textit{89}(12), Artículo 124004. https://doi.org/10.1103/PhysRevD.89.124004

\bibitem[Heydarzade y Vertogradov(2024)]{dinamica_de_photones_bh_cargado} Heydarzade, Y., y Vertogradov, V. (2024). Dynamical photon spheres in charged black holes and naked singularities. \textit{The European Physical Journal C}, \textit{84}(6), Artículo 582. https://doi.org/10.1140/epjc/s10052-024-12935-7

\bibitem[Hobson et al.(2006)]{M_P_Hobson_General_Relativity} Hobson, M. P., Efstathiou, G. P., y Lasenby, A. N. (2006). \textit{General relativity: An introduction for physicists}. Cambridge University Press.

\bibitem[Huang y Mai(2021)]{s10052-021-09180-y} Huang, J. H., y Mai, Z. F. (2021). Superradiance and stability of Kerr-Newman black holes. \textit{The European Physical Journal C}, \textit{81}(5), Artículo 402. https://doi.org/10.1140/epjc/s10052-021-09180-y

\bibitem[Johnson et al.(2024)]{2406.09498v1} Johnson, M. D., Lupsasca, A., Strominger, A., Wong, G. N., Hadar, S., Kapec, D., Narayan, R., Roelofs, F., Alexeev, A., y Galison, P. (2024). Key science goals for the Black Hole Explorer (BHEX). \textit{arXiv preprint arXiv:2406.09498}. https://doi.org/10.48550/arXiv.2406.09498
\bibitem[Johnson et al.(2024)]{Johnson2024} Johnson, M. D., Lupsasca, A., Strominger, A., Wong, G. N., Hadar, S., Kapec, D., Narayan, R., Roelofs, F., Alexeev, A., y Galison, P. (2024). Key science goals for the Black Hole Explorer (BHEX). \textit{arXiv preprint arXiv:2406.09498}. https://doi.org/10.48550/arXiv.2406.09498

\bibitem[Kobialko y Gal'tsov(2024)]{perturbaciones_en_bh} Kobialko, K. V., y Gal'tsov, D. V. (2024). Perturbation theory for gravitational shadows in static spherically symmetric spacetimes. \textit{Physical Review D}, \textit{110}(4), Artículo 044039. https://doi.org/10.1103/PhysRevD.110.044039
\bibitem[Kobialko y Gal'tsov(2024)]{Kobialko2024} Kobialko, K. V., y Gal'tsov, D. V. (2024). Perturbation theory for gravitational shadows in static spherically symmetric spacetimes. \textit{Physical Review D}, \textit{110}(4), Artículo 044039. https://doi.org/10.1103/PhysRevD.110.044039

\bibitem[Li et al.(2025)]{Li2025} Li, X., Chen, Y., y Zhang, J. (2025). Observational signatures of Kerr-Newman black holes with geodesic observers. \textit{The European Physical Journal C}, \textit{85}(1), Artículo 42. https://doi.org/10.1140/epjc/s10052-024-12635-7
\bibitem[Li et al.(2024)]{s10052-024-12635-7} Li, X., Chen, Y., y Zhang, J. (2024). Observational signatures of Kerr-Newman black holes with geodesic observers. \textit{The European Physical Journal C}, \textit{85}(1), Artículo 42. https://doi.org/10.1140/epjc/s10052-024-12635-7

\bibitem[Nayfeh(1993)]{Ali-H.-Nayfeh-Introduction-to-perturbation-techniques-1993-Wiley-VCH} Nayfeh, A. H. (1993). \textit{Introduction to perturbation techniques}. Wiley-VCH.
\bibitem[Nayfeh(1993)]{Nayfeh1993} Nayfeh, A. H. (1993). \textit{Introduction to perturbation techniques}. Wiley-VCH.

\bibitem[Patel et al.(2022)]{Reproducibility_of_the_First_Image_of_a_Black_Hole_in_the_Galaxy_M87_From_the_Event_Horizon_Telescope_Collaboration} Patel, R., Pesce, D. W., Blackburn, L., y Chael, A. (2022). Reproducibility of the first image of a black hole in the galaxy M87 from the Event Horizon Telescope Collaboration. \textit{IEEE Annals of the History of Computing}, \textit{44}(4), 30–42. https://doi.org/10.1109/MAHC.2022.3216854

\bibitem[Raine y Thomas(2015)]{Raine_Thomas_2015} Raine, D., y Thomas, E. (2015). \textit{Black holes: An introduction} (2.ª ed.). Imperial College Press / World Scientific.

\bibitem[Sneppen(2021a)]{Sneppen} Sneppen, A. (2021a). Divergent reflections around the photon sphere of a black hole. \textit{Scientific Reports}, \textit{11}(1), Artículo 14195. https://doi.org/10.1038/s41598-021-93357-w
\bibitem[Sneppen(2021a)]{Sneppen2021} Sneppen, A. (2021a). Divergent reflections around the photon sphere of a black hole. \textit{Scientific Reports}, \textit{11}(1), Artículo 14195. https://doi.org/10.1038/s41598-021-93357-w

\bibitem[Sneppen(2021b)]{pertbuciiones_teoria_en_bh} Sneppen, A. (2021b). Light trajectories and photon ring structure near black holes. \textit{Physical Review D}, \textit{103}(4), Artículo 044031. https://doi.org/10.1103/PhysRevD.103.044031

\bibitem[Téllez Tovar et al.(2025)]{shadows_quintessence_bh} Téllez Tovar, L. O., González, A., y Ramírez, M. (2025). Shadows and photon rings of black holes surrounded by quintessence. \textit{Publicación de Ciencias Básicas Pädi}, \textit{13}(25), 8–15.

\bibitem[Tiede et al.(2022)]{Measuring_Photon_Rings_with_the_ngEHT} Tiede, P., Johnson, M. D., Pesce, D. W., Palumbo, D. C. M., Chang, D. O., y Galison, P. (2022). Measuring photon rings with the ngEHT. \textit{Galaxies}, \textit{10}(6), Artículo 111. https://doi.org/10.3390/galaxies10060111

\bibitem[Valencia Sullca(2021)]{ValenciaSullca2021} Valencia Sullca, J. (2021). \textit{Técnicas de perturbación para ecuaciones diferenciales con incertidumbre y aplicaciones} [Tesis de Licenciatura, Universidad Nacional Mayor de San Marcos]. Repositorio Institucional UNMSM.

\bibitem[Villegas(2022)]{05villegas} Villegas, F. (2022). Sombra del agujero negro de Schwarzschild. \textit{Revista de Investigación de Física}, \textit{25}(3), 40–48. https://doi.org/10.15381/rif.v25i3.22814

\end{thebibliography}

%========================================================
%COMENTAR SI NO DESEA QUE APAREZCA EN EL INDICE GENERAL
%========================================================



%==============================================================
%  SE ADICIONA LAS SECCIONES: ANEXO Y APÉNDICE 
%  COMENTAR SI NO DESEA QUE APAREZCA ESTA PARTE DEL DOCUMENTO
%==============================================================
\appendix{}
\renewcommand{\appendixname}{Anexos}
\renewcommand{\appendixtocname}{Anexos}
\renewcommand{\appendixpagename}{Anexos}
\clearpage
\addappheadtotoc
\appendixpage

%==========================================
%\setcounter{figure}{0}
%\numberwithin{figure}{chapter}
%\numberwithin{table}{section}
%==========================================

%========================= ANEXO 1 ===========================
\chapter{Anexo I: Características de las nubes}
\label{anexoI}

Colocar aquí la información correspondiente.



%========================= ANEXO 2 ===========================
%\chapter{Anexo II: Tablas de doble entrada para el análisis de datos geofísicos}
%\label{anexoII}

%Colocar aquí la información correspondiente. %% Opcional %COMENTAR SI FUERA EL CASO
%==========================================
\appendix{}
\renewcommand{\appendixname}{Apéndices}
\renewcommand{\appendixtocname}{Apéndices}
\renewcommand{\appendixpagename}{Apéndices}
\clearpage
\addappheadtotoc
\appendixpage

%========================= APÉNDICE 1 ===========================
\chapter{Apéndice I: Índice de secas y inundaciones}
\label{apendiceI}

Colocar aquí la información correspondiente.



%========================= APÉNDICE 2 ===========================
%\chapter{Apéndice II: Tablas de doble entrada para el análisis de datos geofísicos}
%\label{apendiceII}

%Colocar aquí la información correspondiente.

 %% Opcional % COMENTAR SI FUERA EL CASO

\end{document}